%% 第9巻 巻末（こたえ・記号カード・変換の地図・定数と単位・公式・チェック表・物理シリーズへ・さくいん・シリーズさくいん・奥付）（第15回編集委員会、組 B9）
%% 第8巻の巻末（ch8-back.tex）と同じ並び・同じ部品。ルビなし。字の色は \bkTT。
%% 委員会の決定9：よく使う公式の節の最初に「変換の地図」（2見開き＝4ページ。偶数ページから：\bkevenpage）。
%%   左1＝前書き・どの表を見る？（流れ図）、右1＝見分け方・表9-2（9-3 の J・J′ の行をふくむ）・表9-4、
%%   左2＝表9-6・表9-7・表9-8、右2＝表9-9・表9-10・表9-11。型は tools/kata9-int.md（頭文字 kixz）。
%%   「こんなとき」のことばと「例題」の番号は、tools/mkmap9.py で本文の表と例題の1行め（場合：…）から写す（手で直さない）。
%%   流れ図「どの表を見る？」は、tools/mkmap9.py --flow で、この巻のすべての例題の場合をたどって確かめた。
%% 委員会の決定13：「つぎの巻へ」のかわりに「物理シリーズへ」。物理シリーズの巻名・章名は phys/series_phys.py から写した。
%%   巻末の最後に、シリーズ全体のさくいん（tex/ch9-seriesidx.tex。tools/mkseriesidx.py で vol1〜9.bkidx から作る）。
%% 記号は委員会の決定どおり：j、s=σ+jω、L と大文字、H(t)、δ(t)（6-7 の δ とは別）、nω0、a0/2、G_2、t_1、矢印のベクトル。
%% 数と式は本文と同じ（scratchpad/back9/check9.py で確かめた）。
\makeatletter
%% ---------------- この巻末だけで使う部品（頭文字 kixz：tools/kata9-int.md の型） ----------------
\colorlet{kixzNc}{bkOrange!85!black}% 名前の枠の線（点線）
\colorlet{kixzSc}{bkTheme}% 実際の操作の枠の線（実線）
\colorlet{kixzSb}{bkTheme!8}% 実際の操作の地
\tikzset{
  kixz dan/.style={font=\sffamily\bfseries\footnotesize, rounded corners=1.2mm, inner xsep=1.4mm, inner ysep=0.4mm, line width=0.9pt},
  kixz dan1/.style={kixz dan, draw=kixzNc, dash pattern=on 2.2pt off 1.4pt, fill=white, text=bkInk},
  kixz dan2/.style={kixz dan, draw=kixzSc, fill=kixzSb, text=bkInk},
  kixz dan3/.style={kixz dan, draw=bkTheme, fill=bkTheme, text=bkOn},
  kixz name/.style={draw=kixzNc, line width=0.9pt, dash pattern=on 2.2pt off 1.4pt, rounded corners=1.2mm, inner xsep=1mm, inner ysep=0.9mm, fill=white},
  kixz seki/.style={draw=kixzSc, line width=0.9pt, rounded corners=1.2mm, inner xsep=1mm, inner ysep=0.9mm, fill=kixzSb},
  kixz tag/.style={font=\sffamily\bfseries\footnotesize, text=bkInk, inner sep=0.4mm},
  kixz q/.style={diamond, aspect=4.6, draw=bkYellow!75!black, line width=0.9pt, fill=bkYellowBg, font=\sffamily\footnotesize, align=flush center, inner sep=0.2mm, text=bkInk},
  kixz do/.style={rectangle, draw=\bkTT, line width=0.9pt, fill=white, rounded corners=0.6mm, font=\sffamily\footnotesize, align=flush center, inner xsep=1.5mm, inner ysep=1.1mm, text=bkInk},
  kixz go/.style={-{Stealth[length=2.2mm]}, line width=0.9pt, bkInk!75, shorten <=1.5mm, shorten >=1.5mm},
  kixz yn/.style={font=\sffamily\bfseries\footnotesize, text=bkInk, fill=white, inner sep=0.5mm},
}
% 実際の操作の4種類の名前（実線の札の字）  1 ふつうの定積分　2 表を引く　3 部分分数　4 行列の計算
\def\kixz@op#1{\ifcase#1\or ふつうの定積分\or 表を引く\or 部分分数\or 行列の計算\fi}
% 段の札  \kixzdan{1|2|3}（名前・直す・楽をする）
\newcommand\kixzdan[1]{\tikz[baseline=(d.base)]\node[kixz dan#1] (d) {\vphantom{あ}\ifcase#1\or 名前\or 直す\or 楽をする\fi};}
% 実線の札（どの道具か）  \kixzop{1|2|3|4}
\newcommand\kixzop[1]{\tikz[baseline=(d.base)]\node[kixz dan2] (d) {\vphantom{あ}\kixz@op{#1}};}
% 場合の札（地図の表の中は小さめ）  \kixzcase{C}
\newcommand\kixzcase[1]{\tikz[baseline=(c.base)]\node[fill=bkTheme, text=bkOn, font=\sffamily\bfseries\footnotesize, rounded corners=0.7mm, inner xsep=0.8mm, inner ysep=0.2mm, minimum width=5mm] (c) {\vphantom{Ag}#1};}
% 表の名札  \kixztab{9-2}
\newcommand\kixztab[1]{\tikz[baseline=(t.base)]\node[fill=bkTheme, text=bkOn, font=\sffamily\bfseries\footnotesize, rounded corners=0.8mm, inner xsep=1.2mm, inner ysep=0.6mm] (t) {\vphantom{Ag}表#1};}
% 式の中の「名前」（点線）と「実際の操作」（実線）。地図では文字の大きさの式（\textstyle）
\newcommand\kixzN[1]{\tikz[baseline=(n.base)]\node[kixz name] (n) {$#1$};}
\newcommand\kixzS[1]{\tikz[baseline=(n.base)]\node[kixz seki] (n) {$#1$};}
\newcommand\kixzNt[1]{\tikz[baseline=(n.base)]{\node[kixz name] (n) {$#1$};\node[kixz tag, anchor=north] at ([yshift=-0.6mm]n.south) {名前};}}
\newcommand\kixzSt[2]{\tikz[baseline=(n.base)]{\node[kixz seki] (n) {$#2$};\node[kixz tag, anchor=north] at ([yshift=-0.6mm]n.south) {\kixz@op{#1}};}}
% 名前と操作の小さな絵（kata9-int の \kixcpic と同じ絵。8 は極と安定、9 は状態方程式）  \kixzpic{1〜9}
\newcommand\kixzpic[1]{\begin{tikzpicture}[x=0.8mm, y=0.8mm, baseline=(current bounding box.center)]
  \ifcase#1\or
  \draw[bkGray, line width=0.4pt] (0,4) -- (14,4);
  \draw[bkInk, line width=0.8pt] (1,4) -- (1,7) -- (7,7) -- (7,1) -- (13,1) -- (13,4);
  \draw[bkTheme, line width=0.8pt, dash pattern=on 1.6pt off 0.9pt, domain=1:13, samples=40, smooth] plot (\x,{4+2.6*cos(30*(\x-1))});
  \draw[bkInk, line width=0.5pt] (1,-0.6) -- (1,-1.6) -- (13,-1.6) -- (13,-0.6);
  \or
  \draw[bkGray, line width=0.4pt] (2,4) -- (12,4) (7,-0.5) -- (7,8.5);
  \draw[bkInk, line width=0.7pt] (7,4) circle (3.6);
  \draw[bkTheme, line width=1.1pt, -{Stealth[length=1.6mm]}] (7,4) -- ({7+3.6*cos(-40)},{4+3.6*sin(-40)});
  \draw[bkInk, line width=0.6pt, -{Stealth[length=1.3mm]}] ({7+5*cos(10)},{4+5*sin(10)}) arc (10:-60:5);
  \or
  \draw[bkInk, line width=0.6pt, {Stealth[length=1.4mm]}-{Stealth[length=1.4mm]}] (0,1) -- (14,1);
  \draw[bkInk, line width=0.8pt] (1.5,1) -- (5.5,1) -- (5.5,6.5) -- (8.5,6.5) -- (8.5,1);
  \draw[bkTheme, line width=0.8pt, dash pattern=on 1.6pt off 0.9pt, domain=1.5:12.5, samples=50, smooth] plot (\x,{1+2.2*(1+cos(80*(\x-7)))/2});
  \or
  \draw[bkInk, line width=0.6pt, -{Stealth[length=1.4mm]}] (1,1) -- (14,1);
  \draw[bkInk, line width=0.6pt] (1,0) -- (1,8.5);
  \fill[bkTheme!25, domain=1:13, samples=40, smooth] (1,1) -- plot (\x,{1+6.5*exp(-0.3*(\x-1))}) -- (13,1) -- cycle;
  \draw[bkTheme, line width=0.8pt, domain=1:13, samples=40, smooth] plot (\x,{1+6.5*exp(-0.3*(\x-1))});
  \or
  \draw[bkGray, line width=0.4pt] (0,4) -- (14,4) (10,-0.5) -- (10,8.5);
  \foreach \kixzx/\kixzy in {3/4,6.5/6.8,6.5/1.2} {\draw[bkRed, line width=1pt] ({\kixzx-1},{\kixzy-1}) -- ({\kixzx+1},{\kixzy+1}) ({\kixzx-1},{\kixzy+1}) -- ({\kixzx+1},{\kixzy-1});}
  \or
  \draw[bkInk, line width=0.6pt, -{Stealth[length=1.4mm]}] (0,4) -- (4,4);
  \draw[bkTheme, line width=0.8pt, fill=bkTheme!12] (4.4,1) rectangle (9.6,7);
  \node[font=\footnotesize, text=bkInk, inner sep=0pt] at (7,4) {$G$};
  \draw[bkInk, line width=0.6pt, -{Stealth[length=1.4mm]}] (10,4) -- (14,4);
  \or
  \draw[bkInk, line width=0.6pt] (1.6,0.5) -- (0.6,0.5) -- (0.6,7.5) -- (1.6,7.5) (5.4,0.5) -- (6.4,0.5) -- (6.4,7.5) -- (5.4,7.5);
  \foreach \kixzx/\kixzy in {2.3/2.3,4.7/2.3,2.3/5.7,4.7/5.7} {\fill[bkInk] (\kixzx,\kixzy) circle (0.5);}
  \draw[bkTheme, line width=0.8pt, domain=0:600, samples=80, smooth] plot ({10.5+3.4*exp(-0.0035*\x)*cos(\x)},{4+3.4*exp(-0.0035*\x)*sin(\x)});
  \or
  \fill[bkGreen!12] (0,-0.5) rectangle (8,8.5);
  \draw[bkGray, line width=0.4pt] (0,4) -- (14,4) (8,-0.5) -- (8,8.5);
  \foreach \kixzx/\kixzy in {3/6.5,3/1.5,11.5/4} {\draw[bkRed, line width=1pt] ({\kixzx-1},{\kixzy-1}) -- ({\kixzx+1},{\kixzy+1}) ({\kixzx-1},{\kixzy+1}) -- ({\kixzx+1},{\kixzy-1});}
  \or
  \draw[bkInk, line width=0.6pt, -{Stealth[length=1.4mm]}] (0,4) -- (3.2,4);
  \draw[bkInk, line width=0.6pt] (4.6,0.5) -- (3.6,0.5) -- (3.6,7.5) -- (4.6,7.5) (9.4,0.5) -- (10.4,0.5) -- (10.4,7.5) -- (9.4,7.5);
  \foreach \kixzx/\kixzy in {5.3/2.3,8.7/2.3,5.3/5.7,8.7/5.7} {\fill[bkInk] (\kixzx,\kixzy) circle (0.5);}
  \draw[bkInk, line width=0.6pt, -{Stealth[length=1.4mm]}] (10.8,4) -- (14,4);
  \fi
\end{tikzpicture}}
% 地図の表の字（9.5pt。行の間をつめる）
\def\kixzfont{\fontsize{9.5}{12.5}\selectfont}
% 地図の表（列：記号／こんなとき→手順／例題。1行めは帯、2行めは見出し）
\NewTblrEnviron{kixzmap}
\SetTblrInner[kixzmap]{colspec={Q[c,m,wd=8.5mm] X[l,m] Q[c,m,wd=17mm]}, width=\linewidth, hlines={0.4pt,bkLine}, rowsep=1pt, colsep=1.5mm, stretch=0,
  cells={font=\kixzfont}, row{1}={abovesep=1.6pt, belowsep=1.6pt}, row{2}={bg=bkVol!14, font=\sffamily\bfseries\footnotesize}}
% 「こんなとき」（太字）と矢印  \kixzif{こんなとき}
\newcommand\kixzto{{\sffamily\bfseries\color{\bkTT}→}}
\newcommand\kixzif[1]{{\sffamily\bfseries #1}\enspace\kixzto\enspace}
% 手順の中で、次の道具へ（部分分数 → 表を引く）
\newcommand\kixzthen{\enspace\kixzto\enspace}
% ほかの章の例題  \kixzo{9-1}{2・4}（例題の列の2行め）
\newcommand\kixzo[2]{{\fontsize{8}{9.5}\selectfont #1：#2}}
% 地図の文（前書きと見分け方）
\def\kixzbody{\fontsize{10}{15}\selectfont}
% 地図の表の帯（1行め。表のはば全体）  \kixzband{中身}
\newcommand\kixzband[1]{\parbox{136mm}{\raggedright\kixzfont\vphantom{表}#1}}
% 地図の小見出し（太字の1行）
\newcommand\kixzhead[1]{\par\addvspace{3pt}\noindent{\sffamily\bfseries\small\color{\bkTT}#1}\par\nobreak\vspace{1.5pt}}
% 巻末の小見出し（第8巻と同じ）
\def\bkbkhead#1{\par\addvspace{4mm}\noindent{\sffamily\bfseries\large\color{\bkTT}#1}\par\nobreak\vspace{2mm}}
% 物理シリーズの巻の色（phys/series_phys.py の VOL_COLORS_CMYK、tex/bk-phys.sty の phVol1〜5 と同じ値。どれも白い字で 4.5:1 以上）
\definecolor{kixzP1}{cmyk}{1,0.60,0,0.10}% 群青（力学）
\definecolor{kixzP2}{cmyk}{0,0.85,0.95,0.18}% 朱赤（熱学）
\definecolor{kixzP3}{cmyk}{0.88,0.05,0.50,0.30}% 青緑（波動）
\definecolor{kixzP4}{cmyk}{0.05,0.45,1,0.35}% 金茶（電磁気学）
\definecolor{kixzP5}{cmyk}{0.65,0.85,0,0.10}% 紫（原子）
% 物理の章の札（物理の巻の色）  \kixzpp{3}{波-4}
\newcommand\kixzpp[2]{\tikz[baseline=(p.base)]\node[fill=kixzP#1, text=white, rounded corners=1mm, inner xsep=1.4mm, inner ysep=0.7mm,
  font=\sffamily\bfseries\fontsize{8}{9}\selectfont] (p) {\vphantom{Ag}#2};}
% 数学の章の札（数学の巻の色。道しるべと同じ）  \kixzmp{9-8}
\newcommand\kixzmp[1]{\bk@pill{#1}}
% 物理シリーズの巻のカード（第8巻の「つぎの巻へ」のカードと同じ形。枠は物理の巻の色）
%   \kixzcard{巻の番号}{巻名}{副題}{どんな「なぜ？」に答えるか}{使う道具}
\def\kixzcardrow{1}
\newcommand\kixzcard[5]{%
  \begin{tcolorbox}[enhanced, nobeforeafter, box align=top, width=0.485\linewidth, colframe=kixzP#1, colback=white, boxrule=0.9pt, arc=2.5mm,
    left=2.5mm, right=2.5mm, top=1.5mm, bottom=2mm, equal height group=kixzcard\kixzcardrow]
  \raggedright\kixzpp{#1}{物理 第#1巻}\enspace{\sffamily\bfseries\small #2}\enspace{\sffamily\footnotesize #3\par}\vspace{0.8mm}
  {\sffamily\bfseries\footnotesize\color{kixzP#1!80!black}#4\par}\vspace{1mm}
  {\footnotesize\setlength{\parindent}{0pt}{\sffamily\bfseries\color{bkInk}使う道具}\enspace #5\par}
  \end{tcolorbox}}
% シリーズさくいんの1行の、巻とページ  \kixzsv{3}{12, 56}（→「第3巻 p.12, 56」）
\newcommand\kixzsv[2]{\mbox{{\sffamily\bfseries 第#1巻}\,p.#2}}
\newcommand\kixzsvsep{、\allowbreak}

%% ---------------- こたえ ----------------
\bkAnswerPages

%% ---------------- 記号カード ----------------
%% 前の巻の記号で、この巻でもよく使うもの（章の番号は、はじめて出た巻の章。この巻の章の \bkkigocard は自動でうしろに集まる）
%% 9-5・9-7・9-10・9-11 のカードは、大きなカードに合わせて \scalebox で縮めている。一覧の記号の列（18 mm）では小さすぎるので、
%% 一覧の中だけ、0.5 以上は 1 に、0.5 より小さいものは 0.75 に読みかえる（章の本文のカードはそのまま。第8巻と同じ方法）
\let\bk@origscalebox\scalebox
\protected\def\scalebox#1{\ifdim#1pt<0.5pt \def\kixz@sb{0.75}\else\ifdim#1pt<1pt \def\kixz@sb{1}\else\def\kixz@sb{#1}\fi\fi
  \expandafter\bk@origscalebox\expandafter{\kixz@sb}}%
\bkKigoPages{%
\bkkigorow{$\Delta$}{デルタ}{変化の量・きざみ。データをとる時間のきざみ $\Delta t$（$f_{\mathrm s}=1/\Delta t$：9-4）、コンピュータの1歩の $\Delta t$（9-10）}{2-3}%
\bkkigorow{\sffamily dB}{デシベル}{倍率を対数で表す単位。振幅の倍率は $20\log_{10}|G|$。直列につなぐと dB のたし算（9-8）}{3-12}%
\bkkigorow{$\omega$}{オメガ}{角周波数〔rad/s〕。周波数の $2\pi$ 倍。$n$ 番めの成分は $n\omega_0$（9-1）}{4-8}%
\bkkigorow{$\lim$}{リミット}{近づく先。最終値の定理 $\lim_{t\to\infty}f(t)=\lim_{s\to0}sF(s)$（9-6）}{5-3}%
\bkkigorow{$\dot{x}$、$\ddot{x}$}{エックス・ドット}{時間で1回・2回微分。状態方程式の $\dot{\vec x}$（9-10）}{5-5}%
\bkkigorow{$\displaystyle\int_a^b$}{インテグラル}{ふつうの定積分。$a_n$・$c_n$・$F(\omega)$・$\mathcal L$ の名前は、これに直して計算する（9-2〜9-5）}{5-10}%
\bkkigorow{$\displaystyle\sum_{k=1}^{n}$}{シグマ}{全部たす。フーリエ級数の和（9-1）、データの和・DFT の和（9-2・9-4）}{5-10}%
\bkkigorow{$e$}{イー}{ネイピア数。$e^{-st}$（9-5）、$e^{-j\omega t}$（9-4）、$e^{At}$（9-11）}{5-14}%
\bkkigorow{$n!$}{エヌのかいじょう}{$\mathcal L[t^n]=n!/s^{n+1}$（9-6）、$e^{At}$ の $(At)^2/2!$（9-11）}{5-17}%
\bkkigorow{$\tau$}{タウ}{時定数〔s〕。1次おくれ $K/(\tau s+1)$ の極は $-1/\tau$（9-8・9-9）。たたみこみの中の時刻は $t_1$ と書く（9-6）}{6-4}%
\bkkigorow{$L$・$C$}{エル・シー}{コイルとコンデンサー。RC・RL・RLC の回路（9-5・9-7・9-10）。$y=C\vec x$ の $C$（はかる行列）とは別}{6-4}%
\bkkigorow{$\omega_{\mathrm n}$}{オメガ・エヌ}{固有角振動数（n は立てた字）。2次おくれの形（9-9）。$n$ 番めの成分 $n\omega_0$ とは別（9-1）}{6-6}%
\bkkigorow{$\zeta$}{ゼータ}{減衰比。2次おくれの行きすぎは $\zeta$ だけで決まる（9-9）}{6-7}%
\bkkigorow{$j$}{ジェイ}{虚数単位。$s=\sigma+j\omega$（9-5）、$e^{jn\omega_0t}$（9-3）。数学・物理の本では $i$}{6-8}%
\bkkigorow{$\bar z$}{ゼット・バー}{共役な複素数。$c_{-n}=\overline{c_n}$（9-3）、共役な極の組（9-7・9-9）}{6-8}%
\bkkigorow{$\arg z$}{アーギュメント}{偏角。周波数応答 $G(j\omega)$ のおくれの角（9-8）}{6-9}%
\bkkigorow{$j\omega$}{ジェイ・オメガ}{$G(s)$ に $s=j\omega$ を入れると周波数応答（9-8）}{6-10}%
\bkkigorow{$s$}{エス}{特性方程式の $s$ と同じもの。ラプラス変換の変数 $s=\sigma+j\omega$（9-5）。極は $s$ の値（9-9）}{6-12}%
\bkkigorow{$\omega_{\mathrm d}$}{オメガ・ディー}{減衰があるときのゆれの角振動数。極 $-\sigma\pm j\omega_{\mathrm d}$ の虚部（9-9）}{6-12}%
\bkkigorow{$X$}{エックス（大文字）}{6-13 の複素振幅は $F_0G(j\omega)$（9-8）。大文字は $s$ の世界の $X(s)$ にも使う（9-5）}{6-13}%
\bkkigorow{$\leftrightarrow$}{たいおうする}{変換の組。$f(t)\leftrightarrow F(s)$（9-5）}{6-15}%
\bkkigorow{$A\vec x$}{エー・エックス}{行列×ベクトル。状態方程式 $\dot{\vec x}=A\vec x+Bu$（9-10）}{7-2}%
\bkkigorow{$I$}{アイ}{単位行列。$(sI-A)^{-1}$（9-8・9-11）、$I+A\,\Delta t$（9-10）}{7-3}%
\bkkigorow{$\det A$}{デターミナント・エー}{行列式。$\det(sI-A)=0$ の解が極（9-9・9-10）}{7-4}%
\bkkigorow{$\lambda$}{ラムダ}{固有値＝極（9-9〜9-11）。9-1 では変数分離の定数にも使う}{7-8}%
\bkkigorow{$D$}{ディー}{対角行列。$e^{At}=Pe^{Dt}P^{-1}$（9-11）}{7-9}%
\bkkigorow{$A^n$}{エー・エヌじょう}{行列を $n$ 回かけたもの。$e^{At}$ の $A^2$・$A^3$（9-11）}{7-10}%
\bkkigorow{$\vec x_n$}{ベクトル・エックス・エヌ}{$n$ 歩めの状態。コンピュータのくり返し $\vec x_{k+1}=A_{\mathrm d}\vec x_k$（9-10）}{7-10}%
}
\let\scalebox\bk@origscalebox

%% 記号カードのつづき：手で書くとき・読むときのちゅうい（各章の「まちがえやすいところ」と、委員会の記号の決定から）
\bkbkhead{記号を書くときのちゅうい}
{\small 手で書くと、ほかの記号とまぎれやすいものと、読みまちがえやすい書き方です。$\mathcal L$・$G(s)$・$e^{At}$ のような記号は、まず「名前」として読みます。}\par\vspace{2mm}
{\small
\begin{longtblr}{colspec={Q[c,m,wd=27mm] Q[c,m,wd=24mm] Q[l,m,wd=75mm]}, rowhead=1, row{1}={bg=bkVol!14, font=\sffamily\bfseries\footnotesize}, hlines={0.4pt,bkLine}, rowsep=0.8pt}
記号 & まぎれやすい & 書くとき・読むときのコツ \\
%% Ver2（全巻の修正事項 No.2）：記号と単位が並ぶときの書き方を1行足した（第3巻の巻末と同じ）
$V$〔V〕 & $V$\,V（VV） & 記号と単位が並ぶときは、単位を〔\ 〕でくくる。「電圧 $V$〔V〕」「時定数 $\tau$〔s〕」。グラフの軸と、点の数の組の単位は（\ ）で「$t$（s）」（9-5・9-8） \\
$j$ & $i$、番号の $j$ & 虚数単位。数学・物理の本の $i$ と同じもの。$\frac{1}{j}=-j$（$\frac1j=j$ ではない）（6-8・9-3） \\
$s=\sigma+j\omega$ & ふつうの数の $s$ & $s$ は複素数。変換表の「$s>a$」は「実部が $a$ より大きい」と読む（9-5・9-6） \\
$\mathcal L[f(t)]$ & $L\times f(t)$ & 筆記体の L は操作の名前で、かけ算ではない。定積分か表に直す（9-5） \\
$X(s)$ & $x(s)$ & 小文字は時間の世界、大文字は $s$ の世界。$sX(s)-x(0)$ のように書き分ける（9-5） \\
$\mathcal L^{-1}$ & $\dfrac{1}{\mathcal L}$ & $-1$ は「逆向き」の印。部分分数に分けて、表を逆に引く（9-7） \\
$\mathcal L[fg]$ & $F(s)\,G(s)$ & ふつうのかけ算の行はない。$F(s)G(s)$ になるのはたたみこみ $f*g$（9-5・9-6） \\
$f*g$ & $f\times g$ & たたみこみ。$e^{-t}*e^{-2t}=e^{-t}-e^{-2t}$。中の時刻は $t_1$（専門書の $\tau$）（9-6） \\
$H(t)$ & $u(t)$、$1(t)$ & ステップ関数。専門書では $u(t)$・$1(t)$ とも書く。この本の $u$ は入れるもの（9-5・9-6） \\
$H(t-t_0)$ & $H(t)-t_0$ & $t_0$ 秒後に0から1になる。おくれた $f(t-t_0)$ には $H(t-t_0)$ をつける（9-6） \\
$\delta(t)$ & 6-7 の $\delta$ & デルタ関数（面積が1の細いパルス）。6-7 の対数減衰率 $\delta$ とは別の使い方（9-6） \\
$n\omega_0$ & $\omega_n$、$\omega_{\mathrm n}$ & $n$ 番めの成分は基本の $n$ 倍。$\omega_{\mathrm n}$（立てた n）は固有角振動数（9-1） \\
$\dfrac{a_0}{2}$ & $a_0$ & 定数の項（平均）は $\frac{a_0}{2}$。$a_n$ の式を $n=0$ にも使えるようにした形（9-1・9-2） \\
$c_n$ の $\dfrac1T$ & $\dfrac2T$ & $a_n$・$b_n$ は $\frac2T$、$c_n$ は $\frac1T$。かけるのは $e^{-jn\omega_0t}$（マイナス）（9-3） \\
$c_{-n}$ & $-c_n$ & 実数の波なら $c_{-n}=\overline{c_n}$（共役）（9-3） \\
$F(\omega)$ の $\omega$ & Hz & 角周波数〔rad/s〕。Hz にするには $2\pi$ でわる（9-4） \\
$f_{\mathrm s}$ & 関数の $f$ & サンプリング周波数〔Hz〕。関数 $f(t)$ と同じページでは、添え字をつけるかことばで（9-4） \\
$\dfrac{2|X_k|}{N}$ & $|X_k|$ & FFT の答えの大きさを振幅にするには $\frac2N$ 倍（9-4） \\
$G(s)$ & 6-13 の $G(\omega)$ & 伝達関数。6-13 の $G(\omega)$ は、分母の大きさの2乗に付けた名前（9-8） \\
$\dfrac{G_1}{1+G_1G_2}$ & $\dfrac{G_1}{1-G_1G_2}$ & 引いてもどすフィードバック。もどす箱は $G_2$（専門書の $H(s)$）（9-8） \\
$20\log_{10}|G|$ & $10\log_{10}|G|$ & 振幅（電圧・ゆれ）の倍率は20倍。電力の比なら10倍（3-12・9-8） \\
$C(sI-A)^{-1}B$ & $BC(sI-A)^{-1}$ & 行列のかける順は変えない（7-3・9-8） \\
×・○ & かけ算・0 & $s$ の平面の極（×）と零点（○）（9-9） \\
$\sigma$ & 標準偏差 & $s=\sigma+j\omega$ の $\sigma$ は実部。9-4 のガウスの形の $\sigma$ は標準偏差（9-4・9-9） \\
$T$ & 周期と温度 & 熱の節の $T$ は温度。周期は使わず、棒の長さ $L$ で $\sin\frac{n\pi x}{L}$（9-1） \\
$u$ & 速度 & この巻の $u$ は入れるもの（$u(t)$・$U(s)$）。物理 力-14 の $u$ は流れの速さ（9-8・9-10） \\
$K$ & ばね・比例の強さ & $u=-K\vec x$ の $K$ は1行の行列。6-14 の結合ばね、比例制御の $K_{\mathrm p}$ とは別（9-10） \\
$A_{\mathrm d}$ の固有値 & $A$ の固有値 & くり返し $\vec x_{k+1}=A_{\mathrm d}\vec x_k$ は大きさが1より小さいか。$\dot{\vec x}=A\vec x$ は実部が負か（9-10） \\
$e^{At}$ & 成分ごとの $e$ & 名前。行列の計算（$Pe^{Dt}P^{-1}$ など）か部分分数に直す（9-11） \\
$Pe^{Dt}P^{-1}$ & $P^{-1}e^{Dt}P$ & $P$ は固有ベクトルを並べた行列。左に $P$、右に $P^{-1}$（9-11） \\
\end{longtblr}\par}

%% 記号カードのつづき：ギリシャ文字（この巻での意味に合わせた）
\par\noindent\begin{minipage}{\linewidth}
\bkbkhead{ギリシャ文字の読み方}
{\small この巻では、時定数 $\tau$、減衰比 $\zeta$、角周波数 $\omega$ など、ギリシャ文字がたくさん出てきます。大文字・小文字・読み方と、この巻での意味です（かっこの中は章）。}\par\vspace{2mm}
\begin{bkcenter}
\footnotesize
\begin{tblr}{colspec={Q[c,m,wd=5mm] Q[c,m,wd=5mm] Q[l,m,wd=17mm] Q[l,m,wd=35mm] Q[c,m,wd=5mm] Q[c,m,wd=5mm] Q[l,m,wd=17mm] Q[l,m,wd=35mm]}, row{1}={bg=bkVol!14, font=\sffamily\bfseries\footnotesize}, hlines={0.4pt,bkLine}, vline{5}={0.8pt,bkVol}, rowsep=0.6pt, colsep=2pt}
大 & 小 & 読み方 & この巻での意味 & 大 & 小 & 読み方 & この巻での意味 \\
$\mathrm{A}$ & $\alpha$ & アルファ & 熱拡散率（9-1・9-3・9-11） & $\mathrm{N}$ & $\nu$ & ニュー & ― \\
$\mathrm{B}$ & $\beta$ & ベータ & 積和の式の角（9-2） & $\Xi$ & $\xi$ & クシー & 専門書の周波数 $\xi$（9-4） \\
$\Gamma$ & $\gamma$ & ガンマ & 送電線の $\gamma\ell$、$M^2=\gamma^2I$（9-11） & $\mathrm{O}$ & $o$ & オミクロン & ― \\
$\Delta$ & $\delta$ & デルタ & $\delta(t)$ はデルタ関数（9-6）。$\Delta t$ はきざみ（9-4・9-10） & $\Pi$ & $\pi$ & パイ & 円周率 \\
$\mathrm{E}$ & $\varepsilon$ & イプシロン & パルスの幅（9-6）、ずれ（9-12） & $\mathrm{P}$ & $\rho$ & ロー & ― \\
$\mathrm{Z}$ & $\zeta$ & ゼータ & 減衰比（9-9） & $\Sigma$ & $\sigma$ & シグマ & $\Sigma$ は和。$\sigma$ は $s$ の実部（9-5・9-9）、ガウスの形のはば（9-4） \\
$\mathrm{H}$ & $\eta$ & イータ & ― & $\mathrm{T}$ & $\tau$ & タウ & 時定数（9-5〜9-9） \\
$\Theta$ & $\theta$ & シータ & 角、棒のかたむき（9-9・9-10）。$\Theta(t)$ は時間だけの関数（9-1） & $\Upsilon$ & $\upsilon$ & ウプシロン & ― \\
$\mathrm{I}$ & $\iota$ & イオタ & ― & $\Phi$ & $\phi$ & ファイ & 専門書の $\Phi(t)$（9-11） \\
$\mathrm{K}$ & $\kappa$ & カッパ & ― & $\mathrm{X}$ & $\chi$ & カイ & ― \\
$\Lambda$ & $\lambda$ & ラムダ & 固有値＝極（9-9〜9-11）、変数分離の定数（9-1） & $\Psi$ & $\psi$ & プサイ & ― \\
$\mathrm{M}$ & $\mu$ & ミュー & 1歩の行列の固有値（9-10）、変数分離の定数（9-1） & $\Omega$ & $\omega$ & オメガ & 角周波数。$\omega_0$（9-1）、$\omega_{\mathrm n}$・$\omega_{\mathrm d}$（9-9） \\
\end{tblr}
\end{bkcenter}
{\footnotesize 大文字が英語のアルファベットと同じ形の文字（A、B、E など）は、記号としてはほとんど使いません。デルタ関数の $\delta(t)$ と 6-7 の対数減衰率 $\delta$ は、同じ文字の別の使い方です。時定数の $\tau$ とまぎれないように、たたみこみの中の時刻は $t_1$ と書きます（専門書では $\tau$）。}
\end{minipage}


%% ---------------- 定数・単位・公式 ----------------
%% 変換の地図は3見開き（6ページ）なので、地図を偶数ページ（左）から始める（\bkevenpage）
\bkbacktitle{よく使う定数・単位・公式}
%% 地図の前のページ（奇数）を、単位の表で埋める（まとめ役の指示：メモのページを残さない）
\bksubsection{1}{この巻で使う単位}
{\footnotesize
\begin{longtblr}{colspec={Q[l,m,wd=25mm] Q[l,m,wd=27mm] Q[l,m,wd=60mm] Q[c,m,wd=10mm]}, rowhead=1, row{1}={bg=bkVol!14, font=\sffamily\bfseries\footnotesize}, hlines={0.4pt,bkLine}, rowsep=1.2pt}
量 & 単位（読み方） & ほかの単位との関係・この巻の例 & 章 \\
時間 & s、ms、分 & $\qty{1}{ms}=\qty{0.001}{s}$。冷めるお湯・エアコンは分で数え、そのときの $s$ の単位は1/分 & {9-0\\9-5} \\
周期 & s & 周期 $T$＝1÷（周波数）。50\,Hz なら \qty{0.02}{s} & 9-1 \\
周波数 & Hz（ヘルツ）、kHz & 1秒にくり返す回数（振動数と同じ）。$\qty{1}{kHz}=\qty{1000}{Hz}$ & 9-0 \\
角周波数 & rad/s & $\omega=2\pi\times$（周波数）。50\,Hz なら $100\pi\approx314$\,rad/s & 9-1 \\
ラプラス変換の $s$ & 1/s（1/分） & $st$ は単位のない数。時間を分で数えると、$s$ の単位は1/分 & 9-5 \\
サンプリング周波数 & Hz & 1秒に値をとる回数 $f_{\mathrm s}=1/\Delta t$。CD は \num{44100}\,Hz & 9-4 \\
倍率の dB & dB（デシベル） & 振幅の倍率の $20\log_{10}$。$\frac{1}{\sqrt2}$ 倍で約 $-3.01$\,dB、$\frac{1}{10}$ 倍で $-20$\,dB & 9-8 \\
電圧・電流・抵抗 & V、A、Ω（kΩ） & 時定数 $RC$：$\qty{20}{k\ohm}\times\qty{100}{\micro F}=\qty{2}{s}$ & {9-5\\9-7} \\
静電容量・インダクタンス & F（µF）、H（mH） & RLC の回路：10\,mH・1\,µF・20\,Ω & {9-7\\9-9} \\
力・ばね・ダンパー & N、N/m、N·s/m & 6-13 のばね：\qty{0.5}{kg}・\qty{8}{N/m}・\qty{0.4}{N.s/m} & 9-8 \\
回る速さ & rad/s、rpm & $\qty{10}{rad/s}\approx\qty{95.5}{rpm}$（rpm は1分に回る回数） & 9-8 \\
温度・温度の差 & ℃、K（ケルビン） & 温度の差は \qty{1}{K} と 1\,℃ が同じ & 9-1 \\
熱拡散率 & m²/s & 8-13 の棒：$\alpha=1\times10^{-4}$\,m²/s & {9-1\\9-3} \\
\end{longtblr}\par}

\bkevenpage
\bksubsection{2}{変換の地図}
\label{bk:9-back:chizu}
{\kixzbody\raggedright $a_n$・$c_n$・$F(\omega)$・$\mathcal L$・$\mathcal L^{-1}$・$G(s)$・$e^{At}$ は、どれも操作の\textbf{名前}です。この3見開きで、①どの表か（流れ図）→②どの場合か（こんなとき）→③どの道具に直すか（札）→④同じ形の例題、の順に引きます。\par}
\vspace{1mm}

%% どの表を見る？（流れ図）。ひし形は問い、四角はすること。行き止まりはない（表9-8・表9-6 の四角は下の問いへもどる）
%% この巻のすべての例題（103題）で、tools/mkmap9.py --flow でたどって確かめた
\kixzhead{どの表を見る？}
\begin{center}
\begin{tikzpicture}[x=1mm, y=0.9mm, kixz q/.append style={aspect=6.2}]
  \node[kixz do, text width=104mm] (s) at (62,0) {記号を読む（名前の段）：何があたえられていて、何がほしいか};
  \node[kixz q] (q1) at (35,-15.4) {名前は $F(\omega)$ か DFT？};
  \node[kixz do, text width=43mm, anchor=west] (r1) at (79,-15.4) {\kixztab{9-4}\enspace フーリエ変換};
  \node[kixz q] (q2) at (35,-30.8) {波を $\sin$・$\cos$ の和で書く？};
  \node[kixz do, text width=43mm, anchor=west] (r2) at (79,-30.8) {\kixztab{9-2}\enspace フーリエ係数};
  \node[kixz q] (q3) at (35,-46.2) {名前は $e^{At}$？};
  \node[kixz do, text width=43mm, anchor=west] (r3) at (79,-46.2) {\kixztab{9-11}\enspace 行列の指数関数};
  \node[kixz q] (q4) at (35,-61.6) {状態方程式を立てる・使う？};
  \node[kixz do, text width=43mm, anchor=west] (r4) at (79,-61.6) {\kixztab{9-10}\enspace 状態方程式};
  \node[kixz q] (q5) at (35,-77) {$G(s)$ を作る・使う？};
  \node[kixz do, text width=43mm, anchor=west] (r5) at (79,-77) {\kixztab{9-8}\enspace 伝達関数\\作ったら、下の問いへ};
  \node[kixz q] (q6) at (35,-92.4) {極・安定かを知りたい？};
  \node[kixz do, text width=43mm, anchor=west] (r6) at (79,-92.4) {\kixztab{9-9}\enspace 極と安定};
  \node[kixz q] (q7) at (35,-107.8) {$t$ の式を $s$ の式に変える？};
  \node[kixz do, text width=43mm, anchor=west] (r7) at (79,-107.8) {\kixztab{9-6}\enspace ラプラス変換\\変えたら、下の問いへ};
  \node[kixz q] (q8) at (35,-123.2) {$s$ の式を $t$ の式にもどす？};
  \node[kixz do, text width=43mm, anchor=west] (r8) at (79,-123.2) {\kixztab{9-7}\enspace 逆ラプラス変換};
  \node[kixz do, text width=62mm] (rz) at (35,-137.5) {$s$ の式・$G(s)$ が答え（もどさない）};
  \node[kixz do, text width=112mm] (end) at (62,-151) {確かめ：着いた表の章の流れ図で場合を決めて直し、その章の確かめへ};
  \draw[kixz go] (s.south -| q1) -- (q1);
  % 「はい」は右へ、「いいえ」は下へ（「いいえ」の札は線の左）
  \foreach \kixzq/\kixzr in {q1/r1,q2/r2,q3/r3,q4/r4,q5/r5,q6/r6,q7/r7,q8/r8} {\draw[kixz go] (\kixzq) -- node[kixz yn, above] {はい} (\kixzr);}
  \foreach \kixzq/\kixzr in {q1/q2,q2/q3,q3/q4,q4/q5,q5/q6,q6/q7,q7/q8,q8/rz} {\draw[kixz go] (\kixzq) -- node[kixz yn, left=1mm] {いいえ} (\kixzr);}
  % 「下の問いへ」の四角は、四角の下から、左の列の次の問いの上へもどる
  \draw[kixz go, shorten >=0.4mm] (r5.south) |- (35.8,-85.7);
  \draw[kixz go, shorten >=0.4mm] (r7.south) |- (35.8,-116.5);
  % 表に着いた四角は、右はしの1本の線（バス）にまとめて「確かめ」へ
  \foreach \kixzn in {r1,r2,r3,r4,r6,r8,rz} {\draw[bkInk!60, line width=0.7pt, shorten <=1.5mm] (\kixzn.east) -- (\kixzn.east -| 127,0);}
  \draw[bkInk!60, line width=0.7pt] (127,-15.4) -- (127,-151);
  \draw[kixz go, shorten <=0pt] (127,-151) -- (end.east);
\end{tikzpicture}
\end{center}
\vspace{-1mm}
{\small\raggedright ひし形は問い、四角はすること。表9-8・表9-6 の四角のあとは、左の列の次の問いへもどる。表9-2 には、9-3 の $c_n$ の行もある。表9-10 の E は表9-9 の G・表9-11 の I と、H は表9-8 の E と同じ場合。\par}
\clearpage

%% ---- 変換の地図（右のページ1）：見分け方・表9-2 ----
\noindent\begin{minipage}{\linewidth}
\kixzhead{名前と、実際の操作の見分け方}
\noindent\begin{minipage}[t]{0.36\linewidth}\vspace{0pt}\centering
\kixzNt{\mathcal L[e^{-2t}]}\ $=$\ \kixzSt{2}{1/(s+2)}\par\vspace{4mm}
\kixzNt{e^{At}}\ $=$\ \kixzSt{4}{Pe^{Dt}P^{-1}}\par\vspace{4mm}
\end{minipage}\hfill
\begin{minipage}[t]{0.62\linewidth}\raggedright\kixzbody\vspace{0pt}
{\sffamily\bfseries 点線の枠＝名前。}まだ数の出ない記号。何に・何をかけて・どこからどこまで・何でわるかを読む。\par
{\sffamily\bfseries 実線の枠＝実際の操作。}4つの道具：\kixzop{1}（第5巻）\kixzop{2}（9-6）\kixzop{3}（3-8）\kixzop{4}（第7巻）\par
{\sffamily\bfseries ぬりつぶし＝}\kixzdan{3}\enspace 積分しない・対称性・数値。\par
どの名前も\enspace\kixzdan{1}\,→\,\kixzdan{3}\,→\,\kixzdan{2}\,→\,確かめ
\end{minipage}\par\vspace{0.5mm}
\noindent{\small\raggedright 例題の番号は、表のある章の例題。ほかの章の例題は「9-1：2」（9-1 の例題2）の形で2行めに。\par}
\end{minipage}
\par\vspace{1.5mm plus 1fill}

%% map:9-2:ch9-02:kixc  表9-2 フーリエ係数（9-2 の表。行の順は本文の表と同じ）
\noindent\begin{kixzmap}{}
\SetCell[c=3]{l} \kixzband{\kixztab{9-2}\enspace{\sffamily\bfseries\color{\bkTT}フーリエ係数}\enspace\kixzpic{1}\,\kixzN{a_n,\ b_n}\enspace\kixzto\enspace \kixzop{1}} & & \\
記号 & こんなとき\enspace→\enspace 手順（頭の札が、どの道具か） & 例題 \\
\kixzcase{I} & \kixzif{式がなく、はかった値だけ}\kixzdan{3}\enspace $(2/N)\sum f_k\sin n\omega_0t_k$ & {9\\\kixzo{9-1}{5}} \\
\kixzcase{A} & \kixzif{$\sin$・$\cos$ の積や累乗の式（$\sin^3t$、$\cos^2t$ など）}\kixzdan{3}\enspace 積和・2倍角の式（4-9）で和に直す & 1 \\
\kixzcase{G} & \kixzif{平均（直流分）$a_0/2$ だけほしい}\kixzdan{3}\enspace 1周期の面積÷$T$ & 7 \\
\kixzcase{E} & \kixzif{棒・弦のように、0から $L$ だけで決まる形}\kixzdan{3}\enspace 半回転して広げ、$b_n=(2/L)\int_0^Lf\sin(n\pi x/L)\,\dd x$ & {4・5\\\kixzo{9-1}{6・7}} \\
\kixzcase{F} & \kixzif{周期 $T$ が $2\pi$ でない（50\,Hz など）}\kixzdan{3}\enspace $\omega_0=2\pi/T$ にして D・B・C へ & {6\\\kixzo{9-1}{4}} \\
\kixzcase{D} & \kixzif{偶関数（たて軸で折り返して重なる）か奇関数（原点で半回転して重なる）}\kixzdan{3}\enspace 偶なら $b_n=0$、奇なら $a_n=0$。半周期の2倍 & 2・3・6 \\
\kixzcase{B} & \kixzif{区間ごとに一定（方形波・パルス）}\kixzop{1}\enspace 区間に分ける。$\cos n\pi=(-1)^n$ & {2・6\\\kixzo{9-1}{2・4・6・7}} \\
\kixzcase{C} & \kixzif{区間ごとに1次式（のこぎり波・三角波）}\kixzop{1}\enspace 部分積分（5-16） & {3・4・5\\\kixzo{9-1}{3}} \\
\kixzcase{H} & \kixzif{区間ごとに $\sin$・$\cos$（整流した波）}\kixzop{1}\enspace 積和の式で和に直して定積分 & 8 \\
\end{kixzmap}
%% /map
\par\vspace{-0.4pt}
%% map:9-2:ch9-03:kixd  表9-2 のつづき：複素数の係数 c_n（9-3 の表。例題の番号は 9-3 のもの）
\noindent\begin{kixzmap}{}
\SetCell[c=3]{l} \kixzband{\kixztab{9-2}\enspace{\sffamily\bfseries\color{\bkTT}（つづき）複素数の係数（9-3 の行）}\enspace\kixzpic{2}\,\kixzN{c_n}\enspace\kixzto\enspace \kixzop{1}} & & \\
記号 & こんなとき\enspace→\enspace 手順（例題は 9-3 の番号） & 例題 \\
\kixzcase{J$'$} & \kixzif{$f$ が $e^{jn\omega_0t}$ や $\cos n\omega_0t$・$\sin n\omega_0t$ の（定数倍の）和で書いてある}\kixzdan{3}\enspace $c_n=(a_n-jb_n)/2$、$c_{-n}=\overline{c_n}$、$c_0=a_0/2$ & 1・2・9 \\
\kixzcase{A} & \kixzif{$\sin$・$\cos$ の積やべき（$\sin^3t$ など）で書ける}\kixzdan{3}\enspace オイラーの公式で展開し、$e^{jnt}$ の前の数を読む & 3 \\
\kixzcase{G} & \kixzif{平均 $c_0$（直流分）だけほしい}\kixzdan{3}\enspace 1周期の面積÷$T$ & 4 \\
\kixzcase{F} & \kixzif{周期が $2\pi$ でない（50\,Hz なら $T=0.02$\,s）}\kixzdan{3}\enspace $\omega_0=2\pi/T$。このあと J へ & 7 \\
\kixzcase{J} & \kixzif{区間ごとに一定、または区間ごとに1次式}\kixzop{1}\enspace $(1/T)\int_0^Tf\,e^{-jn\omega_0t}\,\dd t$。$e^{-jn\pi}=(-1)^n$ & 5・6・7・8 \\
\end{kixzmap}
%% /map
\par
\clearpage

%% ---- 変換の地図（左のページ2）：表9-4・表9-6 ----
%% map:9-4:ch9-04:kixe  表9-4 フーリエ変換
\noindent\begin{kixzmap}{}
\SetCell[c=3]{l} \kixzband{\kixztab{9-4}\enspace{\sffamily\bfseries\color{\bkTT}フーリエ変換}\enspace\kixzpic{3}\,\kixzN{F(\omega)}\enspace\kixzto\enspace \kixzop{1}\enspace\kixzop{2}} & & \\
記号 & こんなとき\enspace→\enspace 手順（頭の札が、どの道具か） & 例題 \\
\kixzcase{G} & \kixzif{式がなく、はかった値だけ（録音・測定したとびとびの値）}\kixzdan{3}\enspace 積分しない。DFT $X_k=\sum x_me^{-j2\pi km/N}$ を FFT で。振幅は $2|X_k|/N$ & 8・9・10 \\
\kixzcase{F} & \kixzif{くり返す波（$\sin$・$\cos$ の和、方形波など）}\kixzdan{3}\enspace 積分しない。9-3 の $c_n$ で、$\pm n\omega_0$ に立つ線のスペクトル & 7 \\
\kixzcase{E} & \kixzif{知っている形を、ずらした・$\cos$ をかけた}\kixzdan{3}\enspace 性質の表を引く。もとの形の変換は、下の行（A〜D）で & 5・6 \\
\kixzcase{A} & \kixzif{0でない所が有限の区間だけ（パルス）}\kixzop{1}\enspace その区間のふつうの定積分。$F(0)$＝面積 & 1・5 \\
\kixzcase{D} & \kixzif{ガウスの形 $e^{-t^2/(2\sigma^2)}$}\kixzop{1}\enspace 肩を平方完成し、ガウス積分 $\sqrt\pi$（8-6） & 4・6 \\
\kixzcase{B} & \kixzif{$t\geqq0$ だけの $e^{-at}$（だんだん消える）}\kixzop{1}\enspace $\int_0^\infty e^{-(a+j\omega)t}\,\dd t=1/(a+j\omega)$（5-14 の広義積分） & 2 \\
\kixzcase{C} & \kixzif{両側の $e^{-a|t|}$}\kixzop{1}\enspace $t<0$ と $t>0$ に分けて B を2回：$2a/(a^2+\omega^2)$ & 3 \\
\end{kixzmap}
%% /map
\par\vspace{0pt plus 1fill}
%% map:9-6:ch9-06:kixg  表9-6 ラプラス変換
\noindent\begin{kixzmap}{}
\SetCell[c=3]{l} \kixzband{\kixztab{9-6}\enspace{\sffamily\bfseries\color{\bkTT}ラプラス変換}\enspace\kixzpic{4}\,\kixzN{\mathcal L[f(t)]}\enspace\kixzto\enspace \kixzop{2}（行は\enspace\kixzop{1}\enspace で1回作る）} & & \\
記号 & こんなとき\enspace→\enspace 手順（頭の札が、どの道具か） & 例題 \\
\kixzcase{A} & \kixzif{表にある形（$1$、$t^n$、$e^{at}$、$\sin\omega t$、$\cos\omega t$、$\delta(t)$）}\kixzop{2}\enspace 表を引く。まだない行は、定義の積分（\kixzop{1}）で1回だけ作る & {1・2・3・4・7・10・11\\\kixzo{9-5}{1・7}} \\
\kixzcase{B} & \kixzif{いくつかの項の和（定数倍は前に出す）}\kixzdan{3}\enspace 線形：項ごとに変換して、たす & {2・7・11\\\kixzo{9-5}{7}} \\
\kixzcase{C} & \kixzif{$e^{-at}$ に、表にある別の関数（$\sin$・$\cos$・$t^n$ など）がかかっている（$e^{-at}$ だけなら A）}\kixzop{2}\enspace $e^{-at}$ をとった形の $F(s)$ を引き、$s$ を $s+a$ に（$s$ のずらし） & 3 \\
\kixzcase{D} & \kixzif{おくれて始まる（$H(t-t_0)$ がかかっている。$t_0>0$）}\kixzop{2}\enspace おくれをとった形の $F(s)$ を引き、$e^{-st_0}$ をかける（時間のずらし） & 4・7 \\
\kixzcase{E} & \kixzif{$f'$・$f''$ のある式（微分方程式）を変換する}\kixzop{2}\enspace $f'\to sF-f(0)$、$f''\to s^2F-sf(0)-f'(0)$、$\int_0^tf\,\dd t_1\to F/s$ & {5・6\\\kixzo{9-5}{2・3・4・5・6}} \\
\kixzcase{F} & \kixzif{区間ごとに式が変わる（パルス・階段）}\kixzdan{3}\enspace ステップ関数 $H(t-t_0)$ の和に直して、B・D へ & 7 \\
\kixzcase{G} & \kixzif{$t^n$ 以外の関数（$\sin$・$\cos$ など）に $t$ がかかっている（$t^n$ だけなら A）}\kixzop{2}\enspace $t$ をとった形の $F(s)$ を引き、$-\dd F/\dd s$ & 8 \\
\kixzcase{H} & \kixzif{表にない形・データだけで、数値で求めたい}\kixzdan{3}\enspace 定義の積分を、表計算の台形で（$e^{-st}$ が十分小さくなる $t$ まで） & 9 \\
\kixzcase{I} & \kixzif{たたみこみ $(f*g)(t)$ を変換する}\kixzop{2}\enspace それぞれを変換して、かけ算 $F(s)\,G(s)$ & 10 \\
\end{kixzmap}
%% /map
\par
\clearpage

%% ---- 変換の地図（右のページ2）：表9-7・表9-8 ----
%% map:9-7:ch9-07:kixh  表9-7 逆ラプラス変換
\noindent\begin{kixzmap}{}
\SetCell[c=3]{l} \kixzband{\kixztab{9-7}\enspace{\sffamily\bfseries\color{\bkTT}逆ラプラス変換}\enspace\kixzpic{5}\,\kixzN{\mathcal L^{-1}[F(s)]}\enspace\kixzto\enspace \kixzop{3}\kixzthen\kixzop{2}} & & \\
記号 & こんなとき\enspace→\enspace 手順（頭の札が、どの道具か） & 例題 \\
\kixzcase{A} & \kixzif{表の形そのもの（定数倍をのぞく）}\kixzop{2}\enspace 定数倍を外に出して、逆に引く & 1・9・11 \\
\kixzcase{B} & \kixzif{分母が、ちがう1次式のかけ算（$s$ だけの因数もふくむ）}\kixzop{3}\enspace $\sum A/(s-p_i)$（かくす方法）\kixzthen\kixzop{2}\enspace $\sum Ae^{p_it}$ & 2・3・10 \\
\kixzcase{C} & \kixzif{同じ1次式が2回以上かかる（重解）}\kixzop{3}\enspace $A/(s+a)+B/(s+a)^2$\kixzthen\kixzop{2}\enspace $te^{-at}$ も & 4・5 \\
\kixzcase{D} & \kixzif{分母が、実数の1次式に分けられない2次式（判別式が負）1つだけ（ほかの因数なし）}\kixzop{3}\enspace 平方完成 $(s+a)^2+\omega^2$\kixzthen\kixzop{2}\enspace $e^{-at}\cos\omega t$・$e^{-at}\sin\omega t$ & 6 \\
\kixzcase{E} & \kixzif{分けられない2次式と、ほかの因数（1次式・もう1つの2次式）がまざる}\kixzop{3}\enspace 2次式の分子は $Bs+C$\kixzthen\kixzop{2}\enspace B と D の形で & 7・8 \\
\kixzcase{F} & \kixzif{$e^{-st_0}$ がかかっている}\kixzdan{3}\enspace はずしてもどし、$t_0$ だけ右へずらす & 10 \\
\kixzcase{G} & \kixzif{分子の次数が、分母の次数と同じか大きい}\kixzdan{3}\enspace 先にわり算。定数の項は $\delta(t)$ & 11 \\
\kixzcase{H} & \kixzif{最後の値（落ち着く値）だけほしい}\kixzdan{3}\enspace $\lim_{s\to0}sF(s)$（落ち着くときだけ） & 12 \\
\kixzcase{I} & \kixzif{3次以上の分母で、手の計算では1次式・2次式に分けきれない}\kixzdan{3}\enspace Python で極と係数を出し、B・D の形で読む & 13 \\
\end{kixzmap}
%% /map
\par\vspace{0pt plus 1fill}
%% map:9-8:ch9-08:kixi  表9-8 伝達関数
\noindent\begin{kixzmap}{}
\SetCell[c=3]{l} \kixzband{\kixztab{9-8}\enspace{\sffamily\bfseries\color{\bkTT}伝達関数}\enspace\kixzpic{6}\,\kixzN{G(s)}\enspace\kixzto\enspace \kixzop{2}\enspace\kixzop{4}} & & \\
記号 & こんなとき\enspace→\enspace 手順（頭の札が、どの道具か） & 例題 \\
\kixzcase{A} & \kixzif{微分方程式があたえられている}\kixzop{2}\enspace はじめの値0で $\dot y\to sY$。$Y/U$ & 1 \\
\kixzcase{B} & \kixzif{箱が1列につながる（直列）}\kixzdan{3}\enspace かけ算 $G_1G_2$ & 2・9 \\
\kixzcase{C} & \kixzif{同じものが2つの箱に入り、出たものをたす（並列）}\kixzdan{3}\enspace たし算 $G_1+G_2$ & 3 \\
\kixzcase{D} & \kixzif{出たものをはかって、入れる側へもどす（フィードバック）}\kixzdan{3}\enspace $Y=G_1(R-G_2Y)$ を解いて $G_1/(1+G_1G_2)$ & 4・5 \\
\kixzcase{E} & \kixzif{状態方程式 $\dot{\vec x}=A\vec x+B u$、$y=C\vec x$ があたえられている}\kixzop{4}\enspace $G=C(sI-A)^{-1}B$ & 6 \\
\kixzcase{F} & \kixzif{インパルス応答・ステップ応答（時間の動き）がほしい}\kixzop{3}\enspace $G$・$G/s$ を分けて\kixzthen\kixzop{2}\enspace 表9-7でもどす & 7 \\
\kixzcase{G} & \kixzif{$\sin$ を入れたときの倍率とおくれ（周波数応答）がほしい}\kixzdan{3}\enspace $s=j\omega$。$|G(j\omega)|$ と角、$20\log_{10}|G|$\,dB & 8・9 \\
\kixzcase{H} & \kixzif{実験のデータしかない（1次おくれとわかっている）}\kixzdan{3}\enspace 最後の値で $K$、63.2\,\%の時刻で $\tau$ & 10 \\
\end{kixzmap}
%% /map
\par
\clearpage

%% ---- 変換の地図（左のページ3）：表9-9・表9-10 ----
%% map:9-9:ch9-09:kixj  表9-9 極と安定
\noindent\begin{kixzmap}{}
\SetCell[c=3]{l} \kixzband{\kixztab{9-9}\enspace{\sffamily\bfseries\color{\bkTT}極と安定}\enspace\kixzpic{8}\,\kixzN{\text{極}}\enspace\kixzto\enspace 解く・判定の式\enspace\kixzop{4}} & & \\
記号 & こんなとき\enspace→\enspace 手順（頭の札が、どの道具か） & 例題 \\
\kixzcase{A} & \kixzif{分母が1次式（$\tau s+1$ など）}\kixzdan{2}\enspace 1次方程式を解く。$\tau s+1$ なら極は $-1/\tau$。$\tau>0$ なら安定 & 2・6・7 \\
\kixzcase{B} & \kixzif{分母（2次以上）が、$(s-p)$ や $(s+\sigma)^2+\omega^2$ だけでできている（$p$・$\sigma$ は0でもよい。$s^2+4^2$ もふくむ。かけ算してあってもよい）}\kixzdan{3}\enspace 解かない。かっこごとに0と置いて、極を読む & 1・4 \\
\kixzcase{C} & \kixzif{分母が2次式（分かれていない）}\kixzdan{2}\enspace 解の公式（3-10）。$\zeta$・$\omega_{\mathrm n}$ も読む。安定かだけなら、係数がすべて正か & 3・5・8 \\
\kixzcase{D} & \kixzif{分母が3次式（分かれていない）}\kixzdan{3}\enspace 安定かだけなら3次の判定：係数がすべて正で $a_2a_1>a_0$ & 9 \\
\kixzcase{E} & \kixzif{分母が4次以上（分かれていない）}\kixzdan{3}\enspace Python の \texttt{np.roots} で極を出し、実部の符号を見る & 10 \\
\kixzcase{F} & \kixzif{フィードバックでつないだ全体の極}\kixzdan{2}\enspace 全体 $G_1/(1+G_1G_2)$ を通分して、分母の式を作る。そのあと A〜E & 6・7・8・9 \\
\kixzcase{G} & \kixzif{状態方程式 $\dot{\vec x}=A\vec x+B\vec u$ で書いてある}\kixzop{4}\enspace $\det(sI-A)$ を作る（極＝固有値）。そのあと A〜E & 5 \\
\end{kixzmap}
%% /map
\par\vspace{0pt plus 1fill}
%% map:9-10:ch9-10:kixk  表9-10 状態方程式
\noindent\begin{kixzmap}{}
\SetCell[c=3]{l} \kixzband{\kixztab{9-10}\enspace{\sffamily\bfseries\color{\bkTT}状態方程式}\enspace\kixzpic{9}\,\kixzN{\dot{\vec x}=A\vec x+Bu}\enspace\kixzto\enspace \kixzop{4}} & & \\
記号 & こんなとき\enspace→\enspace 手順（A〜D は立て方、E〜H は使い方） & 例題 \\
\kixzcase{A} & \kixzif{2階の微分の式が1つ（ばね・車）}\kixzop{4}\enspace 位置と速度を状態に（$x_1=x$、$x_2=\dot x$）。$\ddot x$ について解いて2行めに & 1・4 \\
\kixzcase{B} & \kixzif{1階の式がいくつか（水そう・部屋の温度）}\kixzop{4}\enspace 量をそのまま状態にし、右辺の係数を行列に並べる & 2 \\
\kixzcase{C} & \kixzif{電気回路（コイル・コンデンサー）}\kixzop{4}\enspace $v_C$ と $i$ を状態に。$C\dot v_C=i$、$L\dot i=$（電圧の残り） & 3 \\
\kixzcase{D} & \kixzif{$\sin\theta$・$\cos\theta$ などが入る（倒立振子・振り子）}\kixzdan{3}\enspace $\sin\theta\approx\theta$、$\cos\theta\approx1$（線形化）。そのあと A & 4 \\
\kixzcase{E} & \kixzif{何も入れないで落ち着くかを知りたい}\kixzop{4}\enspace $\det(A-\lambda I)=0$ を解き、固有値の実部がすべて負か & 5 \\
\kixzcase{F} & \kixzif{落ち着くように、入れ方 $u$ を決めたい}\kixzop{4}\enspace $u=-K\vec x$。$A-BK$ の特性方程式の係数を、選んだ固有値の式とくらべる & 6 \\
\kixzcase{G} & \kixzif{コンピュータで $\Delta t$ ごとにくり返す}\kixzop{4}\enspace $I+(A-BK)\,\Delta t$ の固有値の大きさが、すべて1より小さいか & 7 \\
\kixzcase{H} & \kixzif{伝達関数 $G(s)$ がほしい}\kixzop{4}\enspace $C(sI-A)^{-1}B$（逆行列は 7-5）。表9-8 の E と同じ & 8 \\
\end{kixzmap}
%% /map
\par
\clearpage

%% ---- 変換の地図（右のページ3）：表9-11・名前と操作の対応表 ----
%% map:9-11:ch9-11:kixl  表9-11 行列の指数関数
\noindent\begin{kixzmap}{}
\SetCell[c=3]{l} \kixzband{\kixztab{9-11}\enspace{\sffamily\bfseries\color{\bkTT}行列の指数関数}\enspace\kixzpic{7}\,\kixzN{e^{At}}\enspace\kixzto\enspace \kixzop{4}\enspace\kixzop{3}} & & \\
記号 & こんなとき\enspace→\enspace 手順（頭の札が、どの道具か） & 例題 \\
\kixzcase{A} & \kixzif{$A$ が対角（式どうしがつながっていない）}\kixzop{4}\enspace 成分ごとに $e^{\lambda t}$ & 1 \\
\kixzcase{B} & \kixzif{固有値がちがう実数}\kixzop{4}\enspace 対角化 $e^{At}=Pe^{Dt}P^{-1}$ & 2・7 \\
\kixzcase{C} & \kixzif{固有値が複素数 $\sigma\pm j\omega$（ゆれる）}\kixzop{4}\enspace $N=A-\sigma I$ とすると $N^2=-\omega^2I$。$e^{\sigma t}\{(\cos\omega t)I+(\sin\omega t/\omega)N\}$ & 3 \\
\kixzcase{D} & \kixzif{$A^2$ が $O$、または単位行列の何倍か}\kixzdan{3}\enspace $A^2=O$ なら $I+At$。$M^2=\gamma^2I$ なら $(\cosh\gamma t)I+(\sinh\gamma t/\gamma)M$ & 4・5 \\
\kixzcase{E} & \kixzif{固有値が重なる（重解）}\kixzop{3}\enspace $(sI-A)^{-1}$ の成分を部分分数に。$1/(s+a)^2\to te^{-at}$（表9-7 の C） & 6 \\
\kixzcase{F} & \kixzif{どの形でも（2×2 なら楽）。確かめに}\kixzop{3}\enspace $(sI-A)^{-1}$ を逆行列（7-5）で作り、成分ごとに部分分数で表を逆に引く & 7 \\
\kixzcase{G} & \kixzif{入れるもの $u$ がある}\kixzop{3}\enspace $(sI-A)^{-1}(\vec x_0+BU)$ を表9-7でもどす & 8 \\
\kixzcase{H} & \kixzif{コンピュータの数値でよい（3×3以上など）}\kixzdan{3}\enspace \texttt{expm}、または $(I+At/n)^n$ & 9 \\
\kixzcase{I} & \kixzif{落ち着くかどうかだけ}\kixzdan{3}\enspace $e^{At}$ を作らない。固有値の実部がすべて負か & 10 \\
\end{kixzmap}
%% /map
\par
\par\vspace{0pt plus 1fill}
%% 名前と操作の対応表（kata9-int.md の2章。9-0 の予告・9-13 と同じ形。絵は各表の帯にある）
\noindent\begin{minipage}{\linewidth}
\kixzhead{名前と、実際の操作の対応表}
\begin{bkcenter}
\begin{bktable}{colspec={Q[l,m,26mm] Q[l,m,48mm] Q[l,m,42mm] Q[c,m,10mm]}, row{1}={bg=bkVol!14, font=\sffamily\bfseries\footnotesize}, cells={font=\kixzfont}, hlines={0.4pt,bkLine}, rowsep=1.2pt, colsep=1.5mm}
名前（記号） & 名前の意味 & 実際の操作 & 表 \\
$a_n$・$b_n$・$c_n$ & 1周期で、$\cos$・$\sin$・$e^{-jn\omega_0t}$ をかけてたす & \kixzop{1} & 9-2 \\
$F(\omega)$ & $-\infty$〜$\infty$ で $e^{-j\omega t}$ をかけてたす & \kixzop{1}\,\kixzop{2} & 9-4 \\
$\mathcal L[f(t)]$ & $0$〜$\infty$ で $e^{-st}$ をかけてたす & \kixzop{2} & 9-6 \\
$\mathcal L^{-1}[F(s)]$ & ラプラス変換すると $F(s)$ になる $f(t)$ & \kixzop{3}\,\kixzop{2} & 9-7 \\
$G(s)$ & はじめの値を0にして、出てくるもの÷入れるもの & \kixzop{2}\,\kixzop{4} & 9-8 \\
極 & 伝達関数の分母が0になる $s$ & 解く・\kixzop{4} & 9-9 \\
$\dot{\vec x}=A\vec x+Bu$ & 状態の変わる速さを、今の状態と入れるもので表す & \kixzop{4} & 9-10 \\
$e^{At}$ & $I+At+(At)^2/2!+\cdots$ & \kixzop{4}\,\kixzop{3} & 9-11 \\
\end{bktable}
\end{bkcenter}
\end{minipage}
\clearpage

\bksubsection{3}{変換表（9-6 の見開きと同じ）}
\label{bk:9-back:hyou}
{\small 9-5・9-6 で、1回ずつ自分で作った行の表です。表を引く前に、分子の数と文字の置き方をそろえます（$\frac{1}{s^2+9}=\frac13\times\frac{3}{s^2+9}$ なので $\frac13\sin3t$）。逆に引くとき（9-7）も、この表を使います。}\par\vspace{1mm}
\begin{bkcenter}
%% saikei:kansu（tools/mkmap9.py が 9-6 の表をそのまま写す。手で直さない）
\begin{bktable}{colspec={Q[c,m,30mm] Q[c,m,34mm] X[l,m] Q[c,m,15mm]}, row{2-Z}={font=\small}, rowsep=1.6pt, colsep=3pt}
$f(t)$ & $F(s)=\mathcal L[f]$ & 作り方 & $s$ の条件 \\
$\delta(t)$ & $1$ & 細いパルスの先（9-6⑦） & どの $s$ も \\
$1$（$H(t)$） & $\dfrac{1}{s}$ & 定義の積分（9-5） & $s>0$ \\
$t$ & $\dfrac{1}{s^2}$ & 部分積分（9-5） & $s>0$ \\
$t^n$ & $\dfrac{n!}{s^{n+1}}$ & 部分積分のくり返し（9-6②） & $s>0$ \\
$e^{at}$ & $\dfrac{1}{s-a}$ & 定義の積分（9-5） & $s>a$ \\
$t\,e^{at}$ & $\dfrac{1}{(s-a)^2}$ & $s$ のずらし（9-6④） & $s>a$ \\
$\sin\omega t$ & $\dfrac{\omega}{s^2+\omega^2}$ & オイラーの公式（9-6③） & $s>0$ \\
$\cos\omega t$ & $\dfrac{s}{s^2+\omega^2}$ & オイラーの公式（9-6③） & $s>0$ \\
$e^{-at}\sin\omega t$ & $\dfrac{\omega}{(s+a)^2+\omega^2}$ & $s$ のずらし（9-6④） & $s>-a$ \\
$e^{-at}\cos\omega t$ & $\dfrac{s+a}{(s+a)^2+\omega^2}$ & $s$ のずらし（9-6④） & $s>-a$ \\
$t\sin\omega t$ & $\dfrac{2\omega s}{(s^2+\omega^2)^2}$ & $t$ をかける行（9-6⑥） & $s>0$ \\
$t\cos\omega t$ & $\dfrac{s^2-\omega^2}{(s^2+\omega^2)^2}$ & $t$ をかける行（9-6⑥） & $s>0$ \\
$H(t-t_0)$ & $\dfrac{e^{-st_0}}{s}$ & 時間のずらし（9-6⑤） & $s>0$ \\
パルス（高さ1、0〜$t_0$） & $\dfrac{1-e^{-st_0}}{s}$ & ステップの和（9-6⑤） & どの $s$ も \\
\end{bktable}
\bktabcaption{ラプラス変換の表（関数）。$a$、$\omega$、$t_0$ は定数（$\omega>0$、$t_0>0$）。$s$ の条件は、$s$ が複素数なら「実部が」と読む。}
%% /saikei
\end{bkcenter}
\begin{bkcenter}
%% saikei:seishitsu（tools/mkmap9.py が 9-6 の表をそのまま写す。手で直さない）
\begin{bktable}{colspec={Q[l,m,28mm] Q[c,m,34mm] Q[c,m,32mm] X[l,m]}, row{2-Z}={font=\small}, rowsep=4pt, colsep=3pt}
性質 & 時間の世界 & $s$ の世界 & 作り方 \\
線形 & $af+bg$ & $aF+bG$ & 積分の性質（9-5⑨） \\
微分 & $f'$ & $sF-f(0)$ & 部分積分（9-5⑥） \\
2回微分 & $f''$ & $s^2F-sf(0)-f'(0)$ & 微分の行を2回（9-5⑦） \\
積分 & $\displaystyle\int_0^tf\,\dd t_1$ & $\dfrac{F}{s}$ & 微分の行から（9-6⑥） \\
$s$ のずらし & $e^{-at}f(t)$ & $F(s+a)$ & 肩のたし算（9-6④） \\
時間のずらし & $f(t-t_0)H(t-t_0)$ & $e^{-st_0}F(s)$ & 置換積分（9-6⑤） \\
$t$ をかける & $t\,f(t)$ & $-\dfrac{\dd F}{\dd s}$ & $s$ で微分（9-6⑥） \\
たたみこみ & $(f*g)(t)$ & $F(s)\,G(s)$ & 積分の順番を入れかえ（9-6⑧） \\
初期値の定理 & $f(0)$ & $\displaystyle\lim_{s\to\infty}sF(s)$ & 微分の行から（9-6⑨） \\
最終値の定理 & $\displaystyle\lim_{t\to\infty}f(t)$ & $\displaystyle\lim_{s\to0}sF(s)$ & 落ち着くときだけ（9-6⑨） \\
\end{bktable}
\bktabcaption{ラプラス変換の表（性質）。ふつうのかけ算 $fg$ の行はない（$FG$ にはならない）。}
%% /saikei
\end{bkcenter}
\clearpage

\bksubsection{4}{よく使う数}
{\small
\begin{longtblr}{colspec={Q[l,m,wd=40mm] Q[l,m,wd=72mm] Q[c,m,wd=12mm]}, rowhead=1, row{1}={bg=bkVol!14, font=\sffamily\bfseries\footnotesize}, hlines={0.4pt,bkLine}, rowsep=1pt}
名前 & 値 & 章 \\
$2\pi$ & 6.2832（1回転の角。角周波数＝$2\pi\times$周波数） & 9-1 \\
$\dfrac{4}{\pi}$ & 1.2732（高さ1の方形波の基本波の大きさ） & {9-1\\9-3} \\
$\dfrac{\pi}{4}$ & 0.7854（$1-\frac13+\frac15-\frac17+\cdots$。4-9 の4項で0.724） & 9-1 \\
$\dfrac{\pi^2}{8}$ & 1.2337（$1+\frac19+\frac{1}{25}+\cdots$。パーセバルの等式から） & 9-2 \\
ギブズの山 & 高さ1の方形波で約1.179（項をふやしても、とびの約9\,\%の行きすぎが残る） & 9-1 \\
$\dfrac{1}{\sqrt2}$ & 0.707（折れ点の倍率。$-3.01$\,dB） & 9-8 \\
$\sqrt{2\pi}$ & 2.5066（$e^{-t^2/2}$ のフーリエ変換の $F(0)$） & 9-4 \\
$e^{-1}$ & 0.3679。時定数 $\tau$ で、最後の変化の $1-e^{-1}=63.2$\,\%に着く & {9-5\\9-9} \\
1次おくれの着き方 & $\tau$ で63.2\,\%、$2\tau$ で86.5\,\%、$3\tau$ で95.0\,\%、$5\tau$ で99.3\,\% & 9-9 \\
2次おくれの行きすぎ & $\zeta=0.1$ で72.9\,\%、0.3 で37.2\,\%、0.5 で16.3\,\%、0.7 で4.6\,\%、1 で0 & 9-9 \\
$\cosh1$、$\sinh1$ & 1.543、1.175（送電線の行列 $e^{M\ell}$） & 9-11 \\
\end{longtblr}\par}

\Needspace{8\baselineskip}
\bkbkhead{この巻でくり返し使う場面と数}
{\small 同じ場面は、どの章でも同じ数で計算しています。}\par\vspace{1mm}
{\small
\begin{longtblr}{colspec={Q[l,m,wd=21mm] Q[l,m,wd=92mm] Q[c,m,wd=12mm]}, rowhead=1, row{1}={bg=bkVol!14, font=\sffamily\bfseries\footnotesize}, hlines={0.4pt,bkLine}, rowsep=1.6pt}
場面 & 式と数 & 章 \\
方形波 & 4-9 の $\sin\theta+\frac13\sin3\theta+\cdots$（90°で4項0.724、たし続けると $\frac{\pi}{4}$）。高さ1なら $b_n=\frac{4}{n\pi}$（奇数番）、$c_n=\frac{2}{jn\pi}$。ギブズの山1.179、ひずみ率48.3\,\% & {9-1\\〜9-4} \\
ラの音 & 440・880・1320\,Hz、大きさ1・0.5・0.25。$f_{\mathrm s}=8000$\,Hz で1秒とって FFT すると、$\frac{2|X_k|}{N}$ が1・0.5・0.25 & {9-0\\9-4} \\
8-13 の棒 & $L=\qty{0.1}{m}$、$\alpha=1\times10^{-4}$\,m²/s、はじめ100\,℃・両はし0\,℃。1項めは $\frac{400}{\pi}e^{-0.0987t}\sin\frac{\pi x}{L}$。まん中が50\,℃になるのは9.47秒、10秒で47.45\,℃ & {9-1\\9-2\\9-11} \\
ギターの弦 & 5弦（A）：\qty{0.65}{m}、110\,Hz、波の速さ $c=\qty{143}{m/s}$。倍音は 220・330・440\,Hz… & {9-1\\9-2} \\
輪の熱 & 半径 \qty{5}{cm}、$\alpha=1\times10^{-4}$\,m²/s。$c_n$ に $e^{-\alpha n^2t/R^2}$ をかける（$n=\pm1$ の項は1秒に0.04の割合で小さくなり、17.3秒で半分） & 9-3 \\
RC の回路 & 6-4：\qty{10}{V}・\qty{20}{k\ohm}・\qty{100}{\micro F}、$\tau=\qty{2}{s}$。$v=10-10e^{-t/2}$、2秒で6.321\,V、6秒で9.502\,V。3-12 のローパスは折れ点100\,Hz（$\tau=\qty{1.592}{ms}$） & {9-3\\9-5\\9-7\\9-8} \\
車A・車B & 車A $\ddot x+3\dot x+2x=0$ → $2e^{-t}-e^{-2t}$（$x(1)=0.600$）。車B $\ddot x+2\dot x+5x=0$ → $e^{-t}(\cos2t+\frac12\sin2t)$（$x(1)=0.014$）。どちらも $x(0)=1$、$\dot x(0)=0$ & {9-5\\〜9-11} \\
6-13 のばね & \qty{0.5}{kg}・\qty{8}{N/m}・\qty{0.4}{N.s/m}、$G=\frac{1}{0.5s^2+0.4s+8}$。0.8\,N・$\omega=2$ でゆすると、ゆれ \qty{0.132}{m}・おくれ7.6° & {9-8\\9-10} \\
RLC の回路 & 10\,mH・1\,µF・20\,Ω、$\zeta=0.1$。10\,V を入れるとコンデンサーは最大17.29\,V（行きすぎ72.9\,\%） & {9-7\\9-9} \\
冷めるお湯 & 6-4：$\tau=20$ 分、80\,℃・部屋20\,℃。$T=20+60e^{-0.05t}$、10分で56.39\,℃ & {9-0\\9-5} \\
熱電対 & $\tau=\qty{30}{s}$。温度が1秒に $\frac{1}{30}$\,℃ずつ上がると、おくれは $1-e^{-t/30}$\,℃（60秒で0.865\,℃）から1\,℃へ & {9-7\\9-9} \\
モーター & $\frac{10}{0.5s+1}$（1\,V で \qty{10}{rad/s}）。比例制御 $K_{\mathrm p}=1$ で、全体の極 $-22$、ずれ9.1\,\% & {9-8\\9-9} \\
エアコン & 外35\,℃、設定26\,℃、部屋の $RC=60$ 分。比例の強さ $RK_{\mathrm p}=4$ で27.8\,℃（時定数12分）、9 で26.9\,℃（6分） & {9-8\\9-9} \\
サスペンション & 1輪ぶん \qty{250}{kg}・\qty{16000}{N/m}・\qty{2000}{N.s/m}。$\omega_{\mathrm n}=8$\,rad/s、$\zeta=0.5$、極 $-4\pm6.93j$。\qty{5}{cm} の段差で車体は最大 \qty{6.49}{cm} & {9-8\\9-9} \\
倒立振子 & 5-18 の $\ddot\theta=4\theta$（固有値 $\pm2$。0.01\,rad が2秒で0.2731\,rad）。$u=-(8\theta+4\dot\theta)$ で固有値 $-2$・$-2$ & {9-0\\9-10\\9-11} \\
水そう & 6-14 の2つの水そう、はじめ $(20,\ 0)$\,cm。10分後 \qty{11.353}{cm} と \qty{8.647}{cm}（$10\pm10e^{-2}$） & {9-7\\9-11} \\
行列の指数関数 & ばね（$\omega=1$、$t=2\pi$）の $(I+At/n)^n$ の対角：$n=63$ で1.3656、630 で1.0318、6283 で1.0031（7-10） & 9-11 \\
\end{longtblr}\par}

\bksubsection{5}{フーリエ級数・フーリエ変換の公式}
{\small 係数や変換を、どの道具に直して求めるかは、変換の地図（\pageref{bk:9-back:chizu}ページ）にまとめました。ここには、よく使う式と結果を集めます。}\par\vspace{1mm}
{\small
\begin{longtblr}{colspec={Q[l,m,wd=26mm] Q[l,m,wd=86mm] Q[c,m,wd=12mm]}, rowhead=1, row{1}={bg=bkVol!14, font=\sffamily\bfseries\footnotesize}, hlines={0.4pt,bkLine}, rowsep=1.6pt}
公式・きまり & 式 & 章 \\
基本角周波数 & $\omega_0=\dfrac{2\pi}{T}$（周期 $T$）。$n$ 番めの成分の角周波数は $n\omega_0$ & 9-1 \\
フーリエ級数 & $f(t)=\dfrac{a_0}{2}+\displaystyle\sum_{n=1}^{\infty}(a_n\cos n\omega_0t+b_n\sin n\omega_0t)$。とびの所では左右の値の平均 & 9-1 \\
フーリエ係数 & $a_n=\dfrac{2}{T}\displaystyle\int_0^Tf(t)\cos n\omega_0t\,\dd t$、\ $b_n=\dfrac{2}{T}\displaystyle\int_0^Tf(t)\sin n\omega_0t\,\dd t$（$a_n$ は $n=0$ にも使う。1周期ならどこからでもよい） & {9-1\\9-2} \\
直交性 & ちがう番号の $\sin n\omega_0t$・$\cos n\omega_0t$ は、かけて1周期で積分すると0。同じものどうしは $\dfrac{T}{2}$ & 9-2 \\
偶関数・奇関数 & $f(-t)=f(t)$ なら $b_n=0$、$f(-t)=-f(t)$ なら $a_n=0$。残りは半周期の2倍 $\dfrac{4}{T}\displaystyle\int_0^{T/2}$ & 9-2 \\
半区間の展開 & 0から $L$ だけの形：$b_n=\dfrac{2}{L}\displaystyle\int_0^Lf(x)\sin\frac{n\pi x}{L}\,\dd x$ & 9-2 \\
パーセバルの等式 & 2乗の平均＝$\Bigl(\dfrac{a_0}{2}\Bigr)^2+\dfrac12\displaystyle\sum(a_n^2+b_n^2)=\sum|c_n|^2$（全体の電力＝成分ごとの電力の和） & {9-2\\9-3} \\
複素フーリエ級数 & $f(t)=\displaystyle\sum_{n=-\infty}^{\infty}c_ne^{jn\omega_0t}$、\ $c_n=\dfrac{1}{T}\displaystyle\int_0^Tf(t)\,e^{-jn\omega_0t}\,\dd t$ & 9-3 \\
$c_n$ と $a_n$・$b_n$ & $c_n=\dfrac{a_n-jb_n}{2}$、$c_{-n}=\overline{c_n}$、$c_0=\dfrac{a_0}{2}$。振幅は $2|c_n|$ & 9-3 \\
熱方程式・弦 & 棒（両はし0）：$T=\displaystyle\sum b_ne^{-\alpha(n\pi/L)^2t}\sin\frac{n\pi x}{L}$。弦：$\sin\dfrac{n\pi x}{L}\cos\dfrac{n\pi ct}{L}$ の和（周波数 $\dfrac{nc}{2L}$） & {9-1\\9-2} \\
フーリエ変換 & $F(\omega)=\displaystyle\int_{-\infty}^{\infty}f(t)\,e^{-j\omega t}\,\dd t$、もどす式 $f(t)=\dfrac{1}{2\pi}\displaystyle\int_{-\infty}^{\infty}F(\omega)\,e^{j\omega t}\,\dd\omega$ & 9-4 \\
よく使う変換 & パルス（高さ1、$|t|<1$）$\dfrac{2\sin\omega}{\omega}$、$e^{-at}$（$t\geqq0$）$\dfrac{1}{a+j\omega}$、$e^{-a|t|}$ は $\dfrac{2a}{a^2+\omega^2}$、$e^{-t^2/(2\sigma^2)}$ は $\sigma\sqrt{2\pi}\,e^{-\sigma^2\omega^2/2}$ & 9-4 \\
変換の性質 & $f(t-t_0)\to e^{-j\omega t_0}F(\omega)$、$f(t)\cos\omega_ct\to\dfrac{F(\omega-\omega_c)+F(\omega+\omega_c)}{2}$、$f'(t)\to j\omega F(\omega)$ & 9-4 \\
DFT & $X_k=\displaystyle\sum_{m=0}^{N-1}x_m\,e^{-j2\pi km/N}$。周波数のきざみ $\dfrac{1}{N\Delta t}$、振幅は $\dfrac{2|X_k|}{N}$ & 9-4 \\
サンプリング定理 & いちばん高い周波数の2倍より大きい $f_{\mathrm s}$ でとる。$\dfrac{f_{\mathrm s}}{2}$ より上は、ほかの周波数に化ける（エイリアシング） & 9-4 \\
\end{longtblr}\par}

\bksubsection{6}{ラプラス変換・伝達関数・極の公式}
{\small 変換の行と性質は、変換表（\pageref{bk:9-back:hyou}ページ）にあります。}\par\vspace{1mm}
{\small
\begin{longtblr}{colspec={Q[l,m,wd=26mm] Q[l,m,wd=86mm] Q[c,m,wd=12mm]}, rowhead=1, row{1}={bg=bkVol!14, font=\sffamily\bfseries\footnotesize}, hlines={0.4pt,bkLine}, rowsep=1.6pt}
公式・きまり & 式 & 章 \\
ラプラス変換 & $\mathcal L[f(t)]=F(s)=\displaystyle\int_0^{\infty}f(t)\,e^{-st}\,\dd t$。小文字は時間の世界、大文字は $s$ の世界 & 9-5 \\
微分の行 & $\mathcal L[f']=sF(s)-f(0)$、$\mathcal L[f'']=s^2F(s)-sf(0)-f'(0)$ & 9-5 \\
微分方程式の解き方 & ①変換する\quad ②$X(s)=$ の形にわる\quad ③部分分数に分ける\quad ④表を逆に引いて、たす & 9-7 \\
2階の式 & $\ddot x+a_1\dot x+a_0x=u$ なら $(s^2+a_1s+a_0)X=U(s)+(s+a_1)x(0)+\dot x(0)$ & 9-7 \\
たたみこみ & $(f*g)(t)=\displaystyle\int_0^tf(t-t_1)\,g(t_1)\,\dd t_1$ を変換すると $F(s)\,G(s)$ & 9-6 \\
初期値・最終値 & $f(0)=\displaystyle\lim_{s\to\infty}sF(s)$、$\displaystyle\lim_{t\to\infty}f(t)=\lim_{s\to0}sF(s)$（落ち着くときだけ） & 9-6 \\
伝達関数 & $G(s)=\dfrac{Y(s)}{U(s)}$（はじめの値をすべて0にする）。$Y=GU$。$G(0)$ はゆっくり変えたときの倍率 & 9-8 \\
つなぎ方 & 直列 $G_1G_2$、並列 $G_1+G_2$、引いてもどすフィードバック $\dfrac{G_1}{1+G_1G_2}$ & 9-8 \\
応答 & インパルス応答は $G(s)$ を、ステップ応答は $\dfrac{G(s)}{s}$ をもどしたもの & {9-6\\9-8} \\
1次おくれ & $\dfrac{K}{\tau s+1}$：ステップ応答 $K(1-e^{-t/\tau})$、極 $-\dfrac{1}{\tau}$ & {9-8\\9-9} \\
2次おくれ & $\dfrac{\omega_{\mathrm n}^2}{s^2+2\zeta\omega_{\mathrm n}s+\omega_{\mathrm n}^2}$：極 $-\zeta\omega_{\mathrm n}\pm j\omega_{\mathrm n}\sqrt{1-\zeta^2}$、行きすぎ $e^{-\pi\zeta/\sqrt{1-\zeta^2}}$ & 9-9 \\
周波数応答 & $s=j\omega$ を入れた $G(j\omega)$。倍率 $|G(j\omega)|$、おくれは角。$20\log_{10}|G(j\omega)|$\,dB。直列は dB のたし算 & 9-8 \\
極・零点 & 伝達関数の分母$=0$ の $s$ が極（×）、分子$=0$ の $s$ が零点（○）。極 $p$ ごとに動きに $e^{pt}$ が入る & 9-9 \\
安定 & すべての極の実部が負。2次の分母は係数がすべて正。3次 $s^3+a_2s^2+a_1s+a_0$ は係数がすべて正で $a_2a_1>a_0$ & 9-9 \\
\end{longtblr}\par}

\bksubsection{7}{状態方程式と行列の指数関数の公式}
{\small
\begin{longtblr}{colspec={Q[l,m,wd=26mm] Q[l,m,wd=86mm] Q[c,m,wd=12mm]}, rowhead=1, row{1}={bg=bkVol!14, font=\sffamily\bfseries\footnotesize}, hlines={0.4pt,bkLine}, rowsep=1.6pt}
公式・きまり & 式 & 章 \\
状態方程式 & $\dot{\vec x}=A\vec x+Bu$、$y=C\vec x$（$\vec x$ は状態、$u$ は入れるもの、$y$ ははかるもの） & 9-10 \\
落ち着く条件 & 何も入れないとき、$A$ の固有値（$\det(A-\lambda I)=0$ の解）の実部がすべて負 & {9-10\\9-11} \\
状態フィードバック & $u=-K\vec x$ を入れると $\dot{\vec x}=(A-BK)\vec x$。$A-BK$ の固有値を選ぶ & 9-10 \\
コンピュータのくり返し & $\vec x_{k+1}=A_{\mathrm d}\vec x_k$、$A_{\mathrm d}=I+(A-BK)\,\Delta t$。固有値の大きさがすべて1より小さければ落ち着く & 9-10 \\
伝達関数 & $G(s)=C(sI-A)^{-1}B$。極は $\det(sI-A)=0$ の解（＝$A$ の固有値） & {9-8\\9-10} \\
行列の指数関数 & $e^{At}=I+At+\dfrac{(At)^2}{2!}+\cdots=\displaystyle\lim_{n\to\infty}\Bigl(I+A\frac{t}{n}\Bigr)^n$。$\vec x(t)=e^{At}\vec x_0$ & 9-11 \\
対角化 & $e^{At}=Pe^{Dt}P^{-1}$（$P$ は固有ベクトルを並べた行列、$e^{Dt}$ は対角に $e^{\lambda t}$） & 9-11 \\
複素数の固有値 & $\sigma\pm j\omega$ なら、$N=A-\sigma I$ として $e^{At}=e^{\sigma t}\Bigl(\cos\omega t\,I+\dfrac{\sin\omega t}{\omega}N\Bigr)$ & 9-11 \\
ラプラス変換で & $\mathcal L[e^{At}]=(sI-A)^{-1}$。入れるものがあれば $\vec X(s)=(sI-A)^{-1}\bigl(\vec x_0+BU(s)\bigr)$ & 9-11 \\
送電線 & $M^2=\gamma^2I$ なら $e^{M\ell}=\cosh\gamma\ell\,I+\dfrac{\sinh\gamma\ell}{\gamma}M$（四端子定数） & 9-11 \\
指数法則 & $e^{A+B}=e^Ae^B$ は $AB=BA$ のときだけ & 9-11 \\
\end{longtblr}\par}

\bksubsection{8}{表計算・Python の使い方}
{\small 表計算は、A列が時刻（または点の番号・角）、1行めが成分の名前（周波数か番号 $n$）です。上から下へ時間が進みます（2-18・6-2 の型）。}\par\vspace{1mm}
{\small
\begin{longtblr}{colspec={Q[l,m,wd=27mm] Q[l,m,wd=58mm] Q[l,m,wd=30mm] Q[c,m,wd=10mm]}, rowhead=1, row{1}={bg=bkVol!14, font=\sffamily\bfseries\footnotesize}, column{2}={font=\footnotesize}, hlines={0.4pt,bkLine}, rowsep=1.6pt, colsep=3pt}
すること & 書き方（例） & 意味・答え & 章 \\
成分をたす & \texttt{=B\$2*SIN(2*PI()*B\$1*\$A3/1000)} & {A列は ms、1行めは Hz、\\2行めは大きさ} & 9-0 \\
角でかく & \texttt{=B\$2*SIN(B\$1*RADIANS(\$A3))} & {SIN はラジアンで\\受けとる} & 9-1 \\
かけて平均の2倍 & {\texttt{=\$B2*SIN(C\$1*RADIANS(\$A2))}\\\texttt{=2*AVERAGE(C2:C37)}} & {方形波の $b_1$ は\\1.2749（本当は1.2732）} & {9-1\\9-2} \\
ラプラス変換を数値で & {\texttt{=EXP(-2*A2)}、\texttt{=B2*C2}\\\texttt{=(D2+D3)/2*0.1}} & {台形でたす。$s=2$ の\\$\mathcal L[1]$ は0.50167} & {9-5\\9-6} \\
{1歩ずつ進める\\（オイラー法）} & {\texttt{=B2+C2*0.1}、\texttt{=C2+D2*0.1}\\\texttt{=-3*C2-2*B2}（D列の加速度）} & {車A。表を引いた答え\\と比べる（6-2 の型）} & 9-7 \\
FFT（Python） & {\texttt{X = np.fft.rfft(x)}\\\texttt{freq = np.fft.rfftfreq(N, 1/fs)}} & {振幅は\\\texttt{2*np.abs(X)/N}} & {9-0\\9-4} \\
8点の DFT & \texttt{np.fft.fft(x)} & 複素数のまま出る & 9-4 \\
方程式の解（極） & \texttt{np.roots([1, 3, 2])} & {係数は次数の高い順\\答えは \texttt{[-2. -1.]}} & {9-7\\9-9} \\
固有値 & \texttt{np.linalg.eigvals(A)} & 複素数の形で出る & {9-10\\9-11} \\
行列の指数関数 & \texttt{scipy.linalg.expm(A*t)} & $e^{At}$ & 9-11 \\
部分分数・応答 & \texttt{scipy.signal.residue}、\texttt{lsim} & 確かめにだけ使う & {9-7\\9-8} \\
\end{longtblr}\par}
\begin{itemize}[leftmargin=1.4\zw, itemsep=2pt, topsep=4pt, label={\color{\bkTT}●}]
\item \texttt{\$} は、コピーしても動かない印です。\texttt{B\$1} はいつも1行め（成分）を、\texttt{\$A3} はいつもA列（時刻）を読みます（9-0）。
\item 表計算の \texttt{SIN}・\texttt{COS} は、角をラジアンで受けとります。度で入れたときは \texttt{RADIANS} で直します（6-9・9-1）。
\item FFT の答え $X_k$ は、大きさを $\frac{2}{N}$ 倍して振幅にします（0\,Hz の成分だけは $\frac{1}{N}$ 倍）。片側だけ見ているので2倍、$N$ 個たしたので $N$ でわります（9-4）。
\item \texttt{np.roots} などの答えは、手で計算したあとの確かめに使います。部分分数の係数や行列の積を、Python に任せて途中の式を省かないようにします（9-7・9-11）。
\end{itemize}

\Needspace{8\baselineskip}
\bkbkhead{プログラムのひな形（9-0・9-4・9-11）}
{\small 上は、ラの音（440・880・1320\,Hz、大きさ1・0.5・0.25）を FFT で分けるプログラム（9-0・9-4）。下は、車A・車B の極と、車A の $e^{At}$（$t=1$）を求めるプログラム（9-9・9-11）です。}\par\vspace{1mm}
\Needspace{17\baselineskip}
\begin{bkcode}[Python]
import numpy as np
fs = 8000                       \# 1秒に何回はかるか
t = np.arange(fs) / fs          \# 1/8000 秒ごとの 8000 個の時刻
x = (np.sin(2*np.pi*440*t) + 0.5*np.sin(2*np.pi*880*t)
     + 0.25*np.sin(2*np.pi*1320*t))   \# 3つの成分をたした音
X = np.fft.rfft(x)              \# FFT（周波数の世界へ）
freq = np.fft.rfftfreq(len(x), 1/fs)  \# 1つ1つの成分の周波数
A = 2*np.abs(X) / len(x)        \# 成分の大きさ（振幅）にそろえる
print(freq[A > 0.1])            \# 大きさが 0.1 より大きい周波数
print(np.round(A[A > 0.1], 3))  \# その大きさ（小数第3位まで）
\end{bkcode}
\begin{bkcodeout}[実行結果]
[ 440.  880. 1320.]
[1.   0.5  0.25]
\end{bkcodeout}
\vspace{1mm}
\Needspace{11\baselineskip}
\begin{bkcode}[Python]
import numpy as np
from scipy.linalg import expm
print(np.roots([1, 3, 2]))      \# 車A の特性方程式の解（極）
print(np.roots([1, 2, 5]))      \# 車B の特性方程式の解（極）
A = np.array([[0, 1],
              [-2, -3]])        \# 車A の状態の行列（7-8）
np.set\_printoptions(precision=4, suppress=True)  \# 小数第4位まで
print(expm(A*1.0))              \# t = 1 の行列の指数関数
\end{bkcode}
\begin{bkcodeout}[実行結果]
[-2. -1.]
[-1.+2.j -1.-2.j]
[[ 0.6004  0.2325]
 [-0.4651 -0.0972]]
\end{bkcodeout}
\vspace{2mm}

\bksubsection{9}{専門書の書き方との対応}
{\small 大学の本では、同じものを別の記号やことばで書きます。9-1〜9-11 の「専門書をのぞいてみよう」で読みかえたものと、委員会で決めたこの本の書き方を集めました。信号処理・制御の本の記号（$X(\omega)$、フィードバックの $H(s)$、ステップの $u(t)$・$1(t)$、ずれの $e(t)$、PID など）は 9-12 の表（\pageref{bk:9-12:senmon}ページ）に、第1〜9巻の記号の訳は 9-13 の⑰（\pageref{bk:9-13:yaku}ページ）にあります。}\par\vspace{1mm}
{\small
\begin{longtblr}{colspec={Q[l,m,wd=52mm] Q[l,m,wd=58mm] Q[c,m,wd=14mm]}, rowhead=1, row{1}={bg=bkVol!14, font=\sffamily\bfseries\footnotesize}, hlines={0.4pt,bkLine}, rowsep=1.2pt}
専門書の書き方 & この本の書き方・意味 & 章 \\
$i$（数学・物理の本の虚数単位） & $j$（電気の本と同じ） & {9-3\\9-4} \\
$f(t)=a_0+\sum(a_n\cos n\omega_0t+\cdots)$ & 定数の項を平均 $a_0$ とする形。この本の $\dfrac{a_0}{2}$ が、その本の $a_0$ & 9-1 \\
$\hat f(n)$、$\displaystyle\int_{-\pi}^{\pi}$、$\sim$ & $c_n$、1周期のどこからでもよい積分、「この級数を対応させる」 & 9-3 \\
第 $n$ 調波の実効値 $I_n$ & $\sqrt2\,I_n=2|c_n|$（振幅の $\frac{1}{\sqrt2}$） & 9-3 \\
$\mathcal F[f]$、$\hat f(\xi)$、$e^{-2\pi ix\xi}$ & $F(\omega)$。周波数 $\xi$〔Hz〕で書くと、$2\pi$ が肩に入り、前の数がつかない & 9-4 \\
$\dfrac{1}{\sqrt{2\pi}}$ を両方の式に & この本は、もどす式だけに $\dfrac{1}{2\pi}$（電気・制御の本の多数派） & 9-4 \\
$F(j\omega)$ & $F(\omega)$（ラプラス変換の $s$ に $j\omega$ を入れた形として書く本） & 9-4 \\
$\mathcal L\{f(t)\}$、$\mathrm{Re}\,s$ & $\mathcal L[f(t)]$、$s$ の実部 $\sigma$ & 9-5 \\
$u(t-a)$（ステップ関数） & $H(t-t_0)$（この本の $u$ は入れるもの） & 9-6 \\
畳み込み $\displaystyle\int_0^tf(\tau)g(t-\tau)\,d\tau$ & たたみこみ。中の時刻 $\tau$ は、この本の $t_1$ & 9-6 \\
ブロムウィッチ積分、ヘビサイドの展開定理、留数 & 名前の上の式（この本では使わない）、かくす方法、部分分数の係数 & 9-7 \\
閉ループ伝達関数 $W(s)$、一巡伝達関数 & フィードバックでつないだ全体、1周の倍率 $G_1G_2$ & 9-8 \\
$N(s)/D(s)$、根、左半平面、漸近安定 & 分子／分母、解、左半分、安定（やがて0に近づく） & 9-9 \\
フルビッツ（ラウス・フルビッツ）の判別法 & 係数による安定の判定（3次は $a_2a_1>a_0$） & 9-9 \\
$\dot{\symbf x}=A\symbf x+\symbf bu$、$y=\symbf c^{\mathsf T}\symbf x$、$u=-\symbf f^{\mathsf T}\symbf x$ & $\dot{\vec x}=A\vec x+Bu$、$y=C\vec x$、$u=-K\vec x$ & 9-10 \\
可制御性行列 $[\symbf b\ \ A\symbf b]$、極配置 & 可制御（2×2 は $\det(B\ \ AB)\neq0$）、固有値を選ぶ & 9-10 \\
z 変換 & コンピュータのくり返しの式を変換する方法（名前だけ） & 9-10 \\
$\exp(At)$、状態遷移行列 $\Phi(t)$ & $e^{At}$（$t$ 秒進める行列） & 9-11 \\
\end{longtblr}\par}

%% ---------------- チェック表 ----------------
\bkbacktitle{できるようになったことチェック表}
{\small 各章の「ふりかえり」を集めた表です。章を学び終えたら1回め、物理シリーズに進む前に2回めのチェックをしよう。できなかった所は、ページにもどって読みなおします。}\par\vspace{2mm}
%% 自動生成（tools/mkchecklist.py）。各章の「ふりかえり」の項目を集めたもの
\begingroup\footnotesize
\def\bkcb{\tikz[baseline=-0.5ex]\draw[\bkTF, line width=0.7pt, rounded corners=0.5pt] (0,-0.3em) rectangle (0.8em,0.5em);}
\begin{longtblr}{colspec={Q[c,m,wd=9mm] X[l,m] Q[c,m,wd=10mm] Q[c,m,wd=9mm] Q[c,m,wd=9mm]}, width=\linewidth, rowhead=1, colsep=3pt,
  row{1}={bg=bkVol!14, font=\sffamily\bfseries\scriptsize}, hlines={0.4pt,bkLine}, rowsep=0.7pt}
章 & できるようになったこと & ページ & 1回め & 2回め \\
{\sffamily\bfseries\color{\bkTT}9-1} & $\sin$ の波をたして方形波に近づくのを表計算で確かめ、スペクトルを読める。 & \pageref{bk:9-1:spec} & \bkcb & \bkcb \\
 & フーリエ級数の形を書き、方形波・のこぎり波の係数を、ふつうの定積分で求められる。 & \pageref{bk:9-1:kakete} & \bkcb & \bkcb \\
 & 棒の熱と弦のゆれを、変数分離とフーリエ級数で解ける。 & \pageref{bk:9-1:netsu} & \bkcb & \bkcb \\
{\sffamily\bfseries\color{\bkTT}9-2} & ちがう周波数の $\sin$ どうしの積を1周期でたすと0になることを、表計算とグラフで確かめられる。 & \pageref{bk:9-2:kakeru} & \bkcb & \bkcb \\
 & 直交性を使って、係数が「かけて積分し、長さの2乗でわる」で取り出せるわけを説明できる。 & \pageref{bk:9-2:toridasu} & \bkcb & \bkcb \\
 & 場合分けの表9-2 と流れ図で、波の形に合った楽な求め方を選び、係数を計算できる。 & \pageref{bk:9-2:basho} & \bkcb & \bkcb \\
{\sffamily\bfseries\color{\bkTT}9-3} & 反対向きに回る2本の矢印をたすと $\cos$ の波になることを作図で確かめ、「負の周波数」の意味を言える。 & \pageref{bk:9-3:ya} & \bkcb & \bkcb \\
 & $a_n$・$b_n$ から複素フーリエ係数 $c_n=\frac{a_n-jb_n}{2}$ を求め、両側のスペクトルを読める。 & \pageref{bk:9-3:kaki} & \bkcb & \bkcb \\
 & 場合分けの表を使って $c_n$ を $e$ の定積分で求め、回路に入れた周期の波や、輪の形の針金の熱を計算できる。 & \pageref{bk:9-3:basho} & \bkcb & \bkcb \\
{\sffamily\bfseries\color{\bkTT}9-4} & 口笛・声・手をたたく音のスペクトルのちがいを説明でき、短いパルスほどスペクトルが広がることを言える。 & \pageref{bk:9-4:oto} & \bkcb & \bkcb \\
 & フーリエ変換 $F(\omega)$ の記号を読み、方形パルスの変換をふつうの定積分で求め、表計算の和で確かめられる。 & \pageref{bk:9-4:teigi} & \bkcb & \bkcb \\
 & 場合分けの表を使って変換を求め、DFT・FFT の結果を読み、サンプリング定理とエイリアシングを説明できる。 & \pageref{bk:9-4:basho} & \bkcb & \bkcb \\
{\sffamily\bfseries\color{\bkTT}9-5} & 「微分→$s$ をかける」の辞書を使って、冷めるお湯の微分方程式を $s$ の式に書きかえ、表でもどして答えを出せる。 & \pageref{bk:9-5:sekai} & \bkcb & \bkcb \\
 & ラプラス変換の記号を「名前」として読み、辞書（変換表）の行を、定義の積分で1回だけ作れる。 & \pageref{bk:9-5:gyou} & \bkcb & \bkcb \\
 & 微分の行 $sF(s)-f(0)$ を部分積分で導き、場合分けの表9-6（A・B・E）を使って微分方程式を変換できる。 & \pageref{bk:9-5:bibun} & \bkcb & \bkcb \\
{\sffamily\bfseries\color{\bkTT}9-6} & 表計算で $f(t)e^{-st}$ を台形でたし、変換表の値と同じになることを確かめられる。 & \pageref{bk:9-6:menseki} & \bkcb & \bkcb \\
 & $t^n$・$\sin$・$\cos$ と、$e^{-at}$ がかかった形の変換表の行を、部分積分・オイラーの公式・$s$ のずらしで作れる。 & \pageref{bk:9-6:tn} & \bkcb & \bkcb \\
 & 場合分けの表9-6と流れ図を使って、おくれ・パルス・デルタ関数・たたみこみをふくむ関数を変換し、初期値・最終値の定理で確かめられる。 & \pageref{bk:9-6:baai} & \bkcb & \bkcb \\
{\sffamily\bfseries\color{\bkTT}9-7} & 微分方程式を「変換する・わり算で解く・部分分数に分ける・表でもどす」の4つの手順で解く流れを、図とオイラー法の表で説明できる。 & \pageref{bk:9-7:tejun} & \bkcb & \bkcb \\
 & 分母がちがう1次式のかけ算になる式を、部分分数に分けて、表を逆に引いて時間の関数にもどせる。 & \pageref{bk:9-7:kurumaA} & \bkcb & \bkcb \\
 & 場合分けの表9-7を使って、重解・分母が0になる $s$ が複素数の式・入力のある式・おくれて始まる式・最後の値だけの場合も、もどし方を選んで計算できる。 & \pageref{bk:9-7:basho} & \bkcb & \bkcb \\
{\sffamily\bfseries\color{\bkTT}9-8} & 「入れると出てくる」箱を、伝達関数 $G(s)=Y(s)\div U(s)$ で表せることを、ブロック線図と表計算で説明できる。 & \pageref{bk:9-8:hako} & \bkcb & \bkcb \\
 & 微分方程式から伝達関数を作り、箱を直列・並列につないだ伝達関数と、ステップ応答・インパルス応答を求められる。 & \pageref{bk:9-8:tsukuru} & \bkcb & \bkcb \\
 & 場合分けの表9-8を使って、フィードバック・状態方程式・周波数応答（dB）の場合も、伝達関数を求めて使える。 & \pageref{bk:9-8:basho} & \bkcb & \bkcb \\
{\sffamily\bfseries\color{\bkTT}9-9} & 伝達関数の分母が0になる $s$（極）を×で打ち、その場所から動き方を当てられる。 & \pageref{bk:9-9:kyoku} & \bkcb & \bkcb \\
 & 1次おくれ・2次おくれの極から、時定数、行きすぎ、最初の山の時刻を求められる。 & \pageref{bk:9-9:niji} & \bkcb & \bkcb \\
 & すべての極の実部が負かどうかで安定を判定し、フィードバックで動く極、3次の判定、残るずれを計算できる。 & \pageref{bk:9-9:antei} & \bkcb & \bkcb \\
{\sffamily\bfseries\color{\bkTT}9-10} & 手のひらに立てた棒の「かたむき」と「たおれる速さ」を状態として、状態の点の動きの図を読める。 & \pageref{bk:9-10:jotai} & \bkcb & \bkcb \\
 & ばね・水そう・回路・倒立振子の式を状態方程式に書き、$A$ の固有値で落ち着くかを判定できる。 & \pageref{bk:9-10:tateru} & \bkcb & \bkcb \\
 & 状態フィードバックで固有値を選び、コンピュータの中のくり返しで落ち着く条件も判定できる。 & \pageref{bk:9-10:fb} & \bkcb & \bkcb \\
{\sffamily\bfseries\color{\bkTT}9-11} & オイラー法のくり返しを表計算で計算し、行列の指数関数 $e^{At}$ に近づくのを確かめられる。 & \pageref{bk:9-11:kiwami} & \bkcb & \bkcb \\
 & 対角の行列と対角化で $e^{At}$ を計算し、状態方程式を解ける。 & \pageref{bk:9-11:taikakuka} & \bkcb & \bkcb \\
 & 場合分けの表で $e^{At}$ の求め方を選び、固有値で落ち着くかを判定できる。 & \pageref{bk:9-11:basho} & \bkcb & \bkcb \\
{\sffamily\bfseries\color{\bkTT}9-12} & 音楽・写真の圧縮で、周波数の成分に分けて小さな成分をけずっていることを説明できる。 & \pageref{bk:9-12:asshuku} & \bkcb & \bkcb \\
 & 通信・MRI・機械の診断で、FFT と直交性が働いていることを言える。 & \pageref{bk:9-12:ofdm} & \bkcb & \bkcb \\
 & PID 制御の3つの項の役割を、極と最終値の定理のことばで説明できる。 & \pageref{bk:9-12:pid} & \bkcb & \bkcb \\
{\sffamily\bfseries\color{\bkTT}9-13} & 7本の路線の名前と終点を言い、全9巻の路線図でたどれる。 & \pageref{bk:9-13:rosen} & \bkcb & \bkcb \\
 & 路線が合流した5つの所を、式と例で説明できる。 & \pageref{bk:9-13:goryu} & \bkcb & \bkcb \\
 & 専門書の式を、名前の段で読み、記号を訳し、4つの道具で計算できる。 & \pageref{bk:9-13:tejun} & \bkcb & \bkcb \\
\end{longtblr}
\endgroup


%% ---------------- 物理シリーズへ（第9巻は最後の巻なので、つぎの巻の予告カードのかわりに） ----------------
%% 物理シリーズの巻名・副題・ひとこと（catch）・章名は phys/series_phys.py から写した。各巻でよく使う数学の巻は 9-13 の⑳の表と同じ。
%% 第9巻を使う物理の章の一覧は 9-13 の⑳にあるので、ここではくり返さない（まとめ役の取り次ぎ、10月29日）。熱-11 の変数分離は 9-1（第15回の決定1）。
\bkbacktitle{物理シリーズへ}
{\kixzbody 小学1年の「リンゴは何個？」から始まった『物理を理解するための数学』は、この第9巻で終わりです。次は物理シリーズ（全5巻）で、ここまでの道具を使って、機械・熱・波・電気・原子のしくみを読みといていきます。物理の本の道しるべと「思い出そう」には、この本の章の番号（9-8 など）が書いてあります。忘れたら、その章にもどれば大丈夫です。\par}
\vspace{2.5mm}
\noindent\begin{minipage}{\linewidth}
\kixzhead{物理の「なぜ？」と、この本の章}
{\small 物理シリーズの章は、身近な「なぜ？」から始まります。その答えにたどり着くまでに使う、この本の章の例です（物理シリーズの構成案から）。}\par\vspace{1mm}
\begin{bkcenter}
\begin{bktable}{colspec={Q[l,m,14mm] Q[l,m,80mm] Q[l,m,36mm]}, row{1}={bg=bkVol!14, font=\sffamily\bfseries\footnotesize}, row{2-Z}={font=\footnotesize}, hlines={0.4pt,bkLine}, rowsep=1pt, colsep=1.5mm}
物理の章 & なぜ？ & この本の章 \\
\kixzpp{1}{力-4} & ジェットコースターは、最初の山より高く上がれる？ & \kixzmp{4-12}\,\kixzmp{8-4}\,\kixzmp{8-7}\,\kixzmp{8-10} \\
\kixzpp{2}{熱-11} & 二重窓にすると、なぜ部屋が冷えにくい？ & \kixzmp{6-4}\,\kixzmp{8-12}\,\kixzmp{8-13}\,\kixzmp{9-1} \\
\kixzpp{3}{波-2} & 洗濯機は、なぜ回り始めにガタガタゆれることがある？ & \kixzmp{3-12}\,\kixzmp{6-10}\,\kixzmp{6-13}\,\kixzmp{9-8} \\
\kixzpp{3}{波-4} & エアコンは、どうやって部屋をちょうどよい温度に保つ？ & \kixzmp{9-6}\,\kixzmp{9-8}\,\kixzmp{9-9}\,\kixzmp{9-10} \\
\kixzpp{3}{波-11} & 動画の車輪が、逆に回って見えることがあるのはなぜ？ & \kixzmp{3-12}\,\kixzmp{3-17}\,\kixzmp{9-4} \\
\kixzpp{4}{電-9} & 充電器をコンセントからぬいても、ランプがしばらく光っているのはなぜ？ & \kixzmp{6-4}\,\kixzmp{6-12}\,\kixzmp{9-7} \\
\kixzpp{5}{原-5} & 電子は、原子の中のどこにいる？ & \kixzmp{6-12}\,\kixzmp{7-8}\,\kixzmp{9-1}\,\kixzmp{9-3} \\
\kixzpp{5}{原-12} & 大昔の木や骨の年代は、どうやってわかる？ & \kixzmp{3-11}\,\kixzmp{6-4}\,\kixzmp{6-14} \\
\kixzpp{5}{原-15} & 原子炉は、どうやって反応を一定に保つ？ & \kixzmp{6-2}\,\kixzmp{6-4}\,\kixzmp{8-12}\,\kixzmp{9-9} \\
\end{bktable}
\end{bkcenter}
\end{minipage}\par
\vspace{1mm}
{\small\raggedright 物理シリーズの各巻でよく使う数学の巻と、第9巻の道具を使う物理の章12の一覧は、9-13 の⑳「数学の道具を、物理で使う」（\pageref{bk:9-13:butsuri}ページ）にあります。\par}
\vspace{2mm}
\bkfukidashi{delta}{第1巻の「リンゴは何個？」から、ここまで来たね。たし算も、微分も、ラプラス変換も、みんな「わからないことを、わかる形に直す」道具だったよ。次は物理シリーズで、その道具を本物の機械や自然に使ってみよう！}
\clearpage
\vspace{3mm}
\noindent\kixzcard{1}{力学}{動く・回る・ゆがむ・流れる}{「なぜ動く？」「なぜ浮かぶ？」に、力で答えよう}{よく使うのは第8巻・第7巻。運動方程式とオイラー法（6-2：力-2）、トラス・ロボットの腕の行列（7-5〜7-7：力-5・力-7）、重積分・発散（第8巻：力-3・力-6・力-13）。}\hfill
\kixzcard{2}{熱学}{エネルギーを使いこなす}{「なぜ冷える？」「なぜ結露する？」に、熱で答えよう}{よく使うのは第8巻・第3巻・第5巻。止める量を書いた偏微分（8-3：熱-2・熱-10）、時定数（6-4：熱-11）、熱方程式の変数分離（9-1：熱-11）。}\par\vspace{1.2mm}
\def\kixzcardrow{2}\noindent\kixzcard{3}{波動}{ゆれと波を読む}{「なぜ空は青い？」「なぜ音が聞こえる？」に、波で答えよう}{第9巻がいちばん活躍する巻。周波数応答（9-8：波-2）、フィードバック制御（9-6・9-8〜9-10：波-4）、弦の固有振動（9-1〜9-3：波-6）、FFT（9-4：波-10・波-11）。}\hfill
\kixzcard{4}{電磁気学}{場から回路・機械・電波へ}{「なぜモータは回る？」「なぜ電波は届く？」に、電気と磁石で答えよう}{よく使うのは第8巻・第6巻。ガウス・ストークスの定理（8-11：電-2・電-6・電-14）、複素数の交流（6-9・6-10：電-10）、過渡現象をラプラス変換で（9-7：電-9）。}\par\vspace{1.2mm}
\def\kixzcardrow{3}\noindent\kixzcard{5}{原子}{量子から半導体・原子力へ}{「LEDはなぜ光る？」「放射線はどこから出る？」に、原子で答えよう}{よく使うのは第6巻・第4巻・第5巻。短いパルスほど周波数が広がる（9-4：原-3）、変数分離（9-1・9-3：原-5）、極と安定（9-9：原-15）。}\hfill
\begin{tcolorbox}[enhanced, nobeforeafter, box align=top, width=0.485\linewidth, colframe=bkVol9, colback=white, boxrule=0.9pt, arc=2.5mm,
    left=2.5mm, right=2.5mm, top=1.5mm, bottom=2mm, equal height group=kixzcard3]
  \raggedright{\sffamily\bfseries\small 物理シリーズの読み方}\par\vspace{1mm}
  {\footnotesize\setlength{\parindent}{0pt}
  ① 第1巻（力学）から順に読みます。\par
  ② 章のはじめの道しるべにある章の番号（6-2 など）が、この本の章です。\par
  ③ 記号や公式を忘れたら、このあとの「シリーズ全体のさくいん」で、第1〜9巻のどこにあるかを引きます。\par
  ④ ★★★★★ の章も「目でわかる」「手でわかる」まで読めば、まず十分です。\par}
\end{tcolorbox}\par


%% ---------------- さくいん ----------------
\bkbacktitle{さくいん}
{\small 数字はページです。その語がはじめて出てくるページや、くわしく説明しているページをのせています。}\par\vspace{2mm}
\begin{multicols}{2}\small\setlength{\parskip}{0pt}
\directlua{bk.printidx()}
\end{multicols}

%% ---------------- シリーズ全体のさくいん（第1〜9巻） ----------------
%% tex/ch9-seriesidx.tex は tools/mkseriesidx.py が tex/vol1〜vol9.bkidx から作る（手で直さない）
\bkbacktitle{シリーズ全体のさくいん}
{\small 第1巻から第9巻までの、さくいんのことばをまとめました。「第3巻 p.12」は、第3巻の12ページです。同じことばがいくつかの巻に出てくるときは、1行に並べました。第1〜3巻のことばは、ルビと分かち書きの空白をはずして並べました。}\par\vspace{2mm}
\begin{multicols}{2}\footnotesize\setlength{\parskip}{0pt}
\IfFileExists{ch9-seriesidx.tex}{%% シリーズ全体のさくいん（第1〜9巻）。tools/mkseriesidx.py が tex/vol1〜9.bkidx から作った。手で直さない。
\bkidxgroup{あ}
\bkidxline{RLC回路}{\kixzsv{6}{260}}
\bkidxline{圧縮}{\kixzsv{9}{540}}
\bkidxline{圧力}{\kixzsv{3}{193}}
\bkidxline{あまり}{\kixzsv{1}{252}}
\bkidxline{あわせる}{\kixzsv{1}{68}}
\bkidxline{安定}{\kixzsv{9}{405}}
\bkidxline{アンペールの法則}{\kixzsv{8}{462}}
\bkidxline{$e$（ネイピア数）}{\kixzsv{5}{340}}
\bkidxline{いくつ分}{\kixzsv{1}{204}}
\bkidxline{移項}{\kixzsv{2}{301}}
\bkidxline{以上・以下・未満}{\kixzsv{2}{42}}
\bkidxline{1次おくれ}{\kixzsv{9}{397}}
\bkidxline{1次関数}{\kixzsv{3}{135}}
\bkidxline{1次近似}{\kixzsv{5}{179}}
\bkidxline{1次不等式}{\kixzsv{3}{145}}
\bkidxline{一次変換}{\kixzsv{7}{77}}
\bkidxline{1次方程式}{\kixzsv{2}{341}}
\bkidxline{一のくらい}{\kixzsv{1}{38}}
\bkidxline{位置ベクトル}{\kixzsv{4}{386}}
\bkidxline{一般解}{\kixzsv{5}{454}}
\bkidxline{一般角}{\kixzsv{4}{272}}
\bkidxline{入れかえても同じ}{\kixzsv{1}{222}}
\bkidxline{因果関係}{\kixzsv{3}{459}}
\bkidxline{因数分解}{\kixzsv{3}{297}}
\bkidxline{$\int $（インテグラル）}{\kixzsv{5}{242}}
\bkidxline{インパルス応答}{\kixzsv{9}{260}}
\bkidxline{うずなし流れ}{\kixzsv{8}{411}}
\bkidxline{うなり}{\kixzsv{4}{321}}
\bkidxline{運動エネルギー}{\kixzsv{6}{221}}
\bkidxline{運動方程式}{\kixzsv{6}{210}}
\bkidxline{エイリアシング}{\kixzsv{9}{192}}
\bkidxline{SI（国際単位系）}{\kixzsv{2}{204}}
\bkidxline{$x$-$t$ グラフ}{\kixzsv{3}{110}}
\bkidxline{エネルギー}{\kixzsv{5}{291}}
\bkidxline{円}{\kixzsv{1}{327}}
\bkidxline{円周角}{\kixzsv{4}{75}}
\bkidxline{円周率（$\pi $）}{\kixzsv{2}{189}}
\bkidxline{円柱座標}{\kixzsv{4}{507}}
\bkidxline{円の式}{\kixzsv{4}{144}}
\bkidxline{円の面積}{\kixzsv{4}{167}\kixzsvsep \kixzsv{5}{101}}
\bkidxline{オイラー角}{\kixzsv{7}{285}}
\bkidxline{オイラーの公式}{\kixzsv{6}{350}}
\bkidxline{オイラーの等式}{\kixzsv{6}{356}}
\bkidxline{オイラー法}{\kixzsv{5}{177}\kixzsvsep \kixzsv{6}{71}}
\bkidxline{OFDM}{\kixzsv{9}{546}}
\bkidxline{おおきいかず}{\kixzsv{1}{22}}
\bkidxline{オームの法則}{\kixzsv{2}{320}}
\bkidxline{おもさ}{\kixzsv{1}{155}}
\bkidxline{おれ線グラフ}{\kixzsv{1}{307}}
\bkidxline{温度（℃）}{\kixzsv{1}{188}}
\bkidxline{温度計}{\kixzsv{1}{187}}
\bkidxgroup{か}
\bkidxline{解（方程式の）}{\kixzsv{2}{338}}
\bkidxline{回帰直線}{\kixzsv{3}{485}}
\bkidxline{階乗}{\kixzsv{5}{424}}
\bkidxline{外心}{\kixzsv{4}{75}}
\bkidxline{外積}{\kixzsv{4}{475}}
\bkidxline{外接円}{\kixzsv{4}{239}}
\bkidxline{外挿}{\kixzsv{3}{113}}
\bkidxline{回転}{\kixzsv{1}{342}}
\bkidxline{回転（rot）}{\kixzsv{8}{400}}
\bkidxline{回転行列}{\kixzsv{7}{243}}
\bkidxline{回転体}{\kixzsv{4}{175}}
\bkidxline{解の公式}{\kixzsv{3}{305}}
\bkidxline{ガウスの定理}{\kixzsv{8}{445}}
\bkidxline{ガウスの法則}{\kixzsv{8}{451}}
\bkidxline{角}{\kixzsv{1}{342}}
\bkidxline{角運動量}{\kixzsv{4}{482}}
\bkidxline{角振動数}{\kixzsv{5}{378}}
\bkidxline{角速度}{\kixzsv{4}{286}\kixzsvsep \kixzsv{5}{370}}
\bkidxline{確率}{\kixzsv{3}{501}}
\bkidxline{かけ算}{\kixzsv{1}{205}}
\bkidxline{×（かける）}{\kixzsv{1}{204}}
\bkidxline{過減衰}{\kixzsv{6}{436}}
\bkidxline{加減法}{\kixzsv{3}{170}}
\bkidxline{かさ}{\kixzsv{1}{152}}
\bkidxline{華氏}{\kixzsv{2}{280}}
\bkidxline{カスティリアーノの定理}{\kixzsv{8}{96}}
\bkidxline{かずの線}{\kixzsv{1}{96}}
\bkidxline{可制御性}{\kixzsv{9}{467}}
\bkidxline{かぞえる}{\kixzsv{1}{21}}
\bkidxline{かぞえる数}{\kixzsv{2}{19}}
\bkidxline{加速度}{\kixzsv{5}{108}}
\bkidxline{片対数グラフ}{\kixzsv{3}{368}}
\bkidxline{かたむき}{\kixzsv{1}{310}}
\bkidxline{傾き}{\kixzsv{3}{131}\kixzsvsep \kixzsv{5}{18}}
\bkidxline{傾きの場}{\kixzsv{6}{40}}
\bkidxline{カテナリー}{\kixzsv{6}{402}}
\bkidxline{過渡応答}{\kixzsv{6}{20}}
\bkidxline{過渡現象}{\kixzsv{9}{322}}
\bkidxline{加法定理}{\kixzsv{4}{310}}
\bkidxline{かわった量}{\kixzsv{1}{84}}
\bkidxline{換算}{\kixzsv{2}{207}}
\bkidxline{関数}{\kixzsv{3}{14, 79}}
\bkidxline{刻み幅（$\Delta t$）}{\kixzsv{6}{74}}
\bkidxline{期待値}{\kixzsv{3}{509}}
\bkidxline{基底}{\kixzsv{7}{365}}
\bkidxline{ギブズ現象}{\kixzsv{9}{44}}
\bkidxline{基本ベクトル}{\kixzsv{4}{446}}
\bkidxline{逆運動学}{\kixzsv{4}{15, 247}\kixzsvsep \kixzsv{7}{433}}
\bkidxline{逆行列}{\kixzsv{7}{192}}
\bkidxline{逆数}{\kixzsv{2}{88}}
\bkidxline{逆ラプラス変換}{\kixzsv{9}{296}}
\bkidxline{球座標}{\kixzsv{4}{508}}
\bkidxline{球座標（重積分）}{\kixzsv{8}{227}}
\bkidxline{Q値}{\kixzsv{6}{475}}
\bkidxline{球の体積}{\kixzsv{4}{173}}
\bkidxline{球の表面積}{\kixzsv{4}{174}}
\bkidxline{行（表の）}{\kixzsv{1}{220}}
\bkidxline{境界条件}{\kixzsv{6}{113}}
\bkidxline{境界値問題}{\kixzsv{6}{113}}
\bkidxline{共振}{\kixzsv{6}{464}}
\bkidxline{共振曲線}{\kixzsv{6}{464}}
\bkidxline{強制振動}{\kixzsv{6}{461}}
\bkidxline{共分散}{\kixzsv{3}{452}}
\bkidxline{共役な複素数}{\kixzsv{6}{290}}
\bkidxline{行列}{\kixzsv{7}{11, 27}}
\bkidxline{行列式}{\kixzsv{7}{145}}
\bkidxline{行列の指数関数}{\kixzsv{9}{496}}
\bkidxline{極}{\kixzsv{9}{394}}
\bkidxline{極形式}{\kixzsv{6}{320}}
\bkidxline{極限}{\kixzsv{5}{62}}
\bkidxline{極座標}{\kixzsv{4}{79}}
\bkidxline{極大・極小}{\kixzsv{5}{158}}
\bkidxline{虚数単位}{\kixzsv{6}{285}}
\bkidxline{虚部}{\kixzsv{6}{287}}
\bkidxline{キルヒホッフの電圧の法則}{\kixzsv{6}{539}}
\bkidxline{キログラム（kg）}{\kixzsv{1}{155}}
\bkidxline{近似（$\approx $）}{\kixzsv{5}{176}}
\bkidxline{偶関数・奇関数}{\kixzsv{9}{98}}
\bkidxline{空間の座標}{\kixzsv{4}{438}}
\bkidxline{九九}{\kixzsv{1}{207}}
\bkidxline{くらい}{\kixzsv{1}{38, 56}}
\bkidxline{グラム（g）}{\kixzsv{1}{155}}
\bkidxline{くり上がり}{\kixzsv{1}{116}}
\bkidxline{グリーンの定理}{\kixzsv{8}{456}}
\bkidxline{くり下がり}{\kixzsv{1}{117}}
\bkidxline{計算の順番}{\kixzsv{2}{234}}
\bkidxline{係数}{\kixzsv{2}{260}\kixzsvsep \kixzsv{3}{173}}
\bkidxline{係数を比べる}{\kixzsv{3}{233}}
\bkidxline{結合法則}{\kixzsv{7}{113}}
\bkidxline{原始関数}{\kixzsv{5}{268}}
\bkidxline{減衰振動}{\kixzsv{6}{257, 445}}
\bkidxline{懸垂線}{\kixzsv{6}{402}}
\bkidxline{減衰比$\zeta $}{\kixzsv{6}{257}}
\bkidxline{原点}{\kixzsv{3}{51}}
\bkidxline{コイル}{\kixzsv{6}{153}}
\bkidxline{項}{\kixzsv{2}{260}}
\bkidxline{交換法則}{\kixzsv{2}{239}}
\bkidxline{広義積分}{\kixzsv{9}{222}}
\bkidxline{校正}{\kixzsv{3}{495}}
\bkidxline{合成（三角関数）}{\kixzsv{4}{317}}
\bkidxline{合成関数}{\kixzsv{3}{88}\kixzsvsep \kixzsv{5}{133}}
\bkidxline{合成関数の微分}{\kixzsv{5}{134}}
\bkidxline{高速フーリエ変換（FFT）}{\kixzsv{9}{188}}
\bkidxline{交点}{\kixzsv{3}{165}}
\bkidxline{合同}{\kixzsv{4}{30}}
\bkidxline{恒等式}{\kixzsv{3}{230}}
\bkidxline{こう配}{\kixzsv{2}{118}}
\bkidxline{勾配}{\kixzsv{8}{144}}
\bkidxline{公理}{\kixzsv{4}{38}}
\bkidxline{交流}{\kixzsv{5}{362}}
\bkidxline{合流（路線の）}{\kixzsv{9}{562}}
\bkidxline{合力}{\kixzsv{4}{348}}
\bkidxline{午後}{\kixzsv{1}{175}}
\bkidxline{誤差}{\kixzsv{2}{130}\kixzsvsep \kixzsv{5}{182}}
\bkidxline{午前}{\kixzsv{1}{175}}
\bkidxline{弧度法}{\kixzsv{4}{64}}
\bkidxline{こみぐあい}{\kixzsv{2}{66}}
\bkidxline{固有角振動数$\omega _{\mathrm n}$}{\kixzsv{6}{211}}
\bkidxline{固有値}{\kixzsv{7}{318}}
\bkidxline{固有値問題}{\kixzsv{6}{231}}
\bkidxline{固有ベクトル}{\kixzsv{7}{318}}
\bkidxline{コンデンサー}{\kixzsv{6}{35, 149}}
\bkidxgroup{さ}
\bkidxline{最終値の定理}{\kixzsv{9}{264}}
\bkidxline{最小値}{\kixzsv{5}{159}}
\bkidxline{最小二乗法}{\kixzsv{3}{482}}
\bkidxline{最大値}{\kixzsv{5}{159}}
\bkidxline{最大値の原理}{\kixzsv{8}{506}}
\bkidxline{最大値・最小値}{\kixzsv{3}{272}}
\bkidxline{$\sin $（サイン）}{\kixzsv{4}{201}}
\bkidxline{$\sin $・$\cos $ の微分}{\kixzsv{5}{367}}
\bkidxline{錯角}{\kixzsv{4}{33}}
\bkidxline{座標}{\kixzsv{3}{51}}
\bkidxline{座標の変換}{\kixzsv{7}{254}}
\bkidxline{差分法}{\kixzsv{8}{555}}
\bkidxline{サラスの方法}{\kixzsv{7}{154}}
\bkidxline{三角関数}{\kixzsv{4}{274}}
\bkidxline{三角形}{\kixzsv{1}{324}}
\bkidxline{三角形の面積}{\kixzsv{4}{164}}
\bkidxline{三角比}{\kixzsv{4}{201}}
\bkidxline{残差}{\kixzsv{3}{477}}
\bkidxline{3重積分}{\kixzsv{8}{225}}
\bkidxline{三相交流}{\kixzsv{6}{362}}
\bkidxline{散布図}{\kixzsv{3}{445}}
\bkidxline{サンプリング周波数}{\kixzsv{9}{184}}
\bkidxline{三平方の定理}{\kixzsv{4}{133}}
\bkidxline{三平方の定理の逆}{\kixzsv{4}{140}}
\bkidxline{四角形}{\kixzsv{1}{324}}
\bkidxline{四角にならべる}{\kixzsv{1}{205}}
\bkidxline{時間}{\kixzsv{1}{172}}
\bkidxline{式の値}{\kixzsv{2}{275}}
\bkidxline{$\Sigma $（シグマ）}{\kixzsv{5}{236}}
\bkidxline{軸を変えて見る}{\kixzsv{7}{361}}
\bkidxline{次元}{\kixzsv{4}{437}}
\bkidxline{時こく}{\kixzsv{1}{172}}
\bkidxline{仕事}{\kixzsv{4}{411}\kixzsvsep \kixzsv{5}{286}}
\bkidxline{四捨五入}{\kixzsv{2}{128}}
\bkidxline{指数}{\kixzsv{2}{152}}
\bkidxline{次数}{\kixzsv{3}{220}\kixzsvsep \kixzsv{5}{426}}
\bkidxline{指数関数}{\kixzsv{3}{335}}
\bkidxline{指数関数の微分}{\kixzsv{5}{343}}
\bkidxline{指数のきまり}{\kixzsv{3}{345}}
\bkidxline{指数分布}{\kixzsv{6}{157}}
\bkidxline{自然対数$\ln $}{\kixzsv{5}{343}}
\bkidxline{下つき文字}{\kixzsv{3}{57}}
\bkidxline{実効値}{\kixzsv{5}{376}}
\bkidxline{実際の操作}{\kixzsv{9}{26}}
\bkidxline{実部}{\kixzsv{6}{287}}
\bkidxline{時定数$\tau $}{\kixzsv{5}{347}}
\bkidxline{四分位数}{\kixzsv{3}{401}}
\bkidxline{斜辺}{\kixzsv{4}{133}}
\bkidxline{重解}{\kixzsv{6}{445}\kixzsvsep \kixzsv{9}{322}}
\bkidxline{周期}{\kixzsv{4}{276}\kixzsvsep \kixzsv{5}{139}}
\bkidxline{重心}{\kixzsv{4}{387}}
\bkidxline{重積分}{\kixzsv{8}{207}}
\bkidxline{終端速度}{\kixzsv{6}{181}}
\bkidxline{10になるくみ}{\kixzsv{1}{70}}
\bkidxline{十のくらい}{\kixzsv{1}{38}}
\bkidxline{10のまとまり}{\kixzsv{1}{36}}
\bkidxline{10のわく}{\kixzsv{1}{23}}
\bkidxline{周波数}{\kixzsv{9}{12}}
\bkidxline{周波数応答}{\kixzsv{9}{371}}
\bkidxline{重力加速度$g$}{\kixzsv{5}{112}}
\bkidxline{J（ジュール）}{\kixzsv{4}{411}}
\bkidxline{16進数}{\kixzsv{2}{370}}
\bkidxline{縮尺}{\kixzsv{4}{101}}
\bkidxline{順運動学}{\kixzsv{4}{14, 246}}
\bkidxline{循環}{\kixzsv{8}{398}}
\bkidxline{瞬間の速さ}{\kixzsv{5}{24}}
\bkidxline{消去法}{\kixzsv{6}{511}}
\bkidxline{条件分岐}{\kixzsv{2}{412}}
\bkidxline{小数}{\kixzsv{1}{268}}
\bkidxline{小数点}{\kixzsv{1}{268}}
\bkidxline{小数の位}{\kixzsv{2}{23}}
\bkidxline{状態}{\kixzsv{9}{447}}
\bkidxline{状態フィードバック}{\kixzsv{9}{462}}
\bkidxline{状態方程式}{\kixzsv{9}{451}}
\bkidxline{常用対数}{\kixzsv{3}{361}}
\bkidxline{初期位相}{\kixzsv{5}{378}}
\bkidxline{初期条件}{\kixzsv{5}{313}\kixzsvsep \kixzsv{6}{42}}
\bkidxline{振動モード}{\kixzsv{6}{509}}
\bkidxline{ジンバルロック}{\kixzsv{7}{289}}
\bkidxline{振幅}{\kixzsv{5}{378}}
\bkidxline{吸いこみ}{\kixzsv{8}{348}}
\bkidxline{すいの体積}{\kixzsv{4}{172}}
\bkidxline{すうじ}{\kixzsv{1}{23}}
\bkidxline{数値}{\kixzsv{2}{20}}
\bkidxline{数値計算}{\kixzsv{6}{71}}
\bkidxline{数直線}{\kixzsv{1}{96}}
\bkidxline{スカラー}{\kixzsv{4}{344}}
\bkidxline{スカラー場}{\kixzsv{8}{11, 179}}
\bkidxline{ステップ応答}{\kixzsv{9}{371}}
\bkidxline{ステップ関数}{\kixzsv{9}{224}}
\bkidxline{ストークスの定理}{\kixzsv{8}{456}}
\bkidxline{ストップウォッチ}{\kixzsv{1}{174}}
\bkidxline{スペクトル}{\kixzsv{9}{15, 45}}
\bkidxline{正規分布}{\kixzsv{3}{510}}
\bkidxline{正規方程式}{\kixzsv{7}{213}}
\bkidxline{正弦定理}{\kixzsv{4}{239}}
\bkidxline{正射影}{\kixzsv{4}{419}}
\bkidxline{整数}{\kixzsv{2}{42}}
\bkidxline{正則}{\kixzsv{7}{194}}
\bkidxline{正の数}{\kixzsv{2}{40}}
\bkidxline{正の字}{\kixzsv{1}{286}}
\bkidxline{成分}{\kixzsv{4}{376}}
\bkidxline{成分（行列の）}{\kixzsv{7}{28}}
\bkidxline{正方形}{\kixzsv{1}{326}}
\bkidxline{積の微分}{\kixzsv{5}{132}}
\bkidxline{積分するだけで解ける形}{\kixzsv{6}{104}}
\bkidxline{積分定数$C$}{\kixzsv{5}{269}}
\bkidxline{積（行列の）}{\kixzsv{7}{106}}
\bkidxline{接線}{\kixzsv{4}{71}\kixzsvsep \kixzsv{5}{64}}
\bkidxline{接線の式}{\kixzsv{5}{60}}
\bkidxline{絶対値}{\kixzsv{3}{59}}
\bkidxline{絶対値（複素数の）}{\kixzsv{6}{317}}
\bkidxline{接頭語}{\kixzsv{2}{205}}
\bkidxline{接平面}{\kixzsv{8}{104}}
\bkidxline{切片}{\kixzsv{3}{136}}
\bkidxline{セル}{\kixzsv{2}{388}}
\bkidxline{0でわる}{\kixzsv{1}{239}}
\bkidxline{1000（せん）}{\kixzsv{1}{41}}
\bkidxline{漸化式}{\kixzsv{7}{390}}
\bkidxline{線形}{\kixzsv{5}{449}\kixzsvsep \kixzsv{9}{231}}
\bkidxline{線形化}{\kixzsv{5}{431}\kixzsvsep \kixzsv{9}{455}}
\bkidxline{線形写像}{\kixzsv{7}{77}}
\bkidxline{線積分}{\kixzsv{8}{256}}
\bkidxline{センチメートル（cm）}{\kixzsv{1}{134}}
\bkidxline{千のくらい}{\kixzsv{1}{41}}
\bkidxline{全微分}{\kixzsv{8}{109}}
\bkidxline{専門書の読み方}{\kixzsv{9}{574}}
\bkidxline{相関}{\kixzsv{3}{446}}
\bkidxline{相関係数}{\kixzsv{3}{455}}
\bkidxline{増加・減少}{\kixzsv{5}{157}}
\bkidxline{双曲線}{\kixzsv{3}{198}\kixzsvsep \kixzsv{6}{393}}
\bkidxline{双曲線関数}{\kixzsv{6}{391}}
\bkidxline{増減表}{\kixzsv{5}{157}}
\bkidxline{相似}{\kixzsv{4}{101}}
\bkidxline{相似（式の形）}{\kixzsv{6}{540}}
\bkidxline{相似比}{\kixzsv{4}{102}}
\bkidxline{相対速度}{\kixzsv{4}{358}}
\bkidxline{速度}{\kixzsv{5}{37, 107}}
\bkidxline{速度を先に進めるオイラー法}{\kixzsv{6}{87}}
\bkidxline{素数}{\kixzsv{2}{80}}
\bkidxgroup{た}
\bkidxline{対角化}{\kixzsv{7}{364}}
\bkidxline{対角行列}{\kixzsv{7}{360}}
\bkidxline{対数}{\kixzsv{3}{361}\kixzsvsep \kixzsv{5}{334}}
\bkidxline{対数減衰率$\delta $}{\kixzsv{6}{255}}
\bkidxline{体積}{\kixzsv{2}{153}}
\bkidxline{だいたい}{\kixzsv{1}{189}}
\bkidxline{代入}{\kixzsv{2}{275}}
\bkidxline{代入法}{\kixzsv{3}{167}}
\bkidxline{ダイヤグラム}{\kixzsv{3}{124}}
\bkidxline{多項式}{\kixzsv{5}{90}}
\bkidxline{たしかめのしき}{\kixzsv{1}{254}}
\bkidxline{たし算}{\kixzsv{1}{68}}
\bkidxline{＋（たす）}{\kixzsv{1}{68}}
\bkidxline{たたみこみ}{\kixzsv{9}{263}}
\bkidxline{たんい}{\kixzsv{1}{134}}
\bkidxline{単位}{\kixzsv{2}{20}}
\bkidxline{単位円}{\kixzsv{4}{273}}
\bkidxline{単位円（固有値の）}{\kixzsv{7}{404}}
\bkidxline{単位行列}{\kixzsv{7}{43, 113}}
\bkidxline{単位のチェック}{\kixzsv{2}{213}}
\bkidxline{単位ベクトル}{\kixzsv{4}{385}}
\bkidxline{単振動}{\kixzsv{5}{371}}
\bkidxline{弾性エネルギー$\frac {1}{2}kx^2$}{\kixzsv{5}{291}}
\bkidxline{ダンパー}{\kixzsv{6}{248}}
\bkidxline{断面二次モーメント}{\kixzsv{8}{240}}
\bkidxline{ちがい}{\kixzsv{1}{83}}
\bkidxline{力の分解}{\kixzsv{4}{209, 380}}
\bkidxline{力のモーメント}{\kixzsv{4}{470}}
\bkidxline{置換積分}{\kixzsv{5}{397}}
\bkidxline{中央値}{\kixzsv{3}{398}}
\bkidxline{鳥かん図}{\kixzsv{8}{34}}
\bkidxline{ちょう点}{\kixzsv{1}{324}}
\bkidxline{頂点}{\kixzsv{3}{262}}
\bkidxline{長方形}{\kixzsv{1}{326}}
\bkidxline{直線}{\kixzsv{1}{324}}
\bkidxline{直角}{\kixzsv{1}{325}}
\bkidxline{直角三角形}{\kixzsv{1}{326}}
\bkidxline{直径}{\kixzsv{1}{327}}
\bkidxline{直交}{\kixzsv{4}{418}}
\bkidxline{直交（関数の）}{\kixzsv{9}{89}}
\bkidxline{直交行列}{\kixzsv{7}{252}}
\bkidxline{通分}{\kixzsv{2}{81}}
\bkidxline{次の瞬間を決めるルール}{\kixzsv{6}{10}}
\bkidxline{底}{\kixzsv{3}{361}}
\bkidxline{定義}{\kixzsv{4}{38}}
\bkidxline{抵抗の係数}{\kixzsv{6}{180}}
\bkidxline{定数}{\kixzsv{3}{27}\kixzsvsep \kixzsv{5}{87}}
\bkidxline{定積分}{\kixzsv{5}{242}}
\bkidxline{テイラー展開}{\kixzsv{5}{425}}
\bkidxline{定理}{\kixzsv{4}{39}}
\bkidxline{データロガー}{\kixzsv{3}{125}}
\bkidxline{デシベル}{\kixzsv{3}{370}}
\bkidxline{デシリットル（dL）}{\kixzsv{1}{152}}
\bkidxline{$\Delta $（デルタ）}{\kixzsv{2}{64}}
\bkidxline{デルタ関数}{\kixzsv{9}{260}}
\bkidxline{展開}{\kixzsv{3}{222}}
\bkidxline{展開図}{\kixzsv{4}{174}}
\bkidxline{電気振動（LC回路）}{\kixzsv{6}{227}}
\bkidxline{天気予報}{\kixzsv{6}{13}}
\bkidxline{電磁波}{\kixzsv{8}{578}}
\bkidxline{伝達関数}{\kixzsv{9}{350}}
\bkidxline{転置}{\kixzsv{7}{40}}
\bkidxline{てんびん}{\kixzsv{1}{154}}
\bkidxline{電力}{\kixzsv{5}{150}}
\bkidxline{電力量}{\kixzsv{5}{292}}
\bkidxline{度（°）}{\kixzsv{1}{343}}
\bkidxline{導関数}{\kixzsv{5}{67}}
\bkidxline{等高線}{\kixzsv{8}{34}}
\bkidxline{等式}{\kixzsv{2}{296}}
\bkidxline{等式の性質}{\kixzsv{2}{298}}
\bkidxline{同次座標}{\kixzsv{7}{290}}
\bkidxline{同次変換行列}{\kixzsv{7}{291}}
\bkidxline{等速直線運動}{\kixzsv{4}{389}}
\bkidxline{等値面}{\kixzsv{8}{46}}
\bkidxline{等比級数}{\kixzsv{7}{400}}
\bkidxline{等比数列}{\kixzsv{7}{399}}
\bkidxline{倒立振子}{\kixzsv{9}{454}}
\bkidxline{同類項}{\kixzsv{2}{260}}
\bkidxline{（〜について）解く}{\kixzsv{2}{300}}
\bkidxline{特異姿勢}{\kixzsv{7}{160}}
\bkidxline{特殊解}{\kixzsv{5}{454}}
\bkidxline{特性インピーダンス}{\kixzsv{7}{441}}
\bkidxline{特性方程式}{\kixzsv{6}{426}\kixzsvsep \kixzsv{7}{330}}
\bkidxline{度数分布表}{\kixzsv{3}{394}}
\bkidxline{ド・モアブルの定理}{\kixzsv{6}{325}}
\bkidxgroup{な}
\bkidxline{内積}{\kixzsv{4}{415}}
\bkidxline{内挿}{\kixzsv{3}{112}}
\bkidxline{ながさ}{\kixzsv{1}{134}}
\bkidxline{ナビエ・ストークス方程式}{\kixzsv{8}{579}}
\bkidxline{ナブラ}{\kixzsv{8}{153}}
\bkidxline{何十・何百のけいさん}{\kixzsv{1}{115}}
\bkidxline{なんしゅうとあといくつ}{\kixzsv{1}{255}}
\bkidxline{2回微分}{\kixzsv{5}{90}}
\bkidxline{2次おくれ}{\kixzsv{9}{400}}
\bkidxline{2次関数}{\kixzsv{3}{269}}
\bkidxline{2次近似}{\kixzsv{5}{422}}
\bkidxline{2次方程式}{\kixzsv{3}{293}}
\bkidxline{2乗に比例}{\kixzsv{3}{259}}
\bkidxline{2進数}{\kixzsv{2}{366}}
\bkidxline{2点の距離}{\kixzsv{4}{142}}
\bkidxline{2倍角}{\kixzsv{4}{313}\kixzsvsep \kixzsv{5}{375}}
\bkidxline{2つの変数の関数（2変数関数）}{\kixzsv{8}{42}}
\bkidxline{ニュートン}{\kixzsv{5}{76}}
\bkidxline{ニュートンの冷却の法則}{\kixzsv{6}{34}}
\bkidxline{音色}{\kixzsv{9}{16}}
\bkidxline{熱方程式}{\kixzsv{8}{543}}
\bkidxline{熱容量$C_{\mathrm {th}}$}{\kixzsv{6}{107}}
\bkidxline{熱容量}{\kixzsv{6}{145}}
\bkidxline{のこり}{\kixzsv{1}{82}}
\bkidxgroup{は}
\bkidxline{＝（は）}{\kixzsv{1}{68}}
\bkidxline{パーセバルの等式}{\kixzsv{9}{103}}
\bkidxline{倍音}{\kixzsv{9}{16, 65}}
\bkidxline{バイト}{\kixzsv{2}{369}}
\bkidxline{ハイパボリック・タンジェント（tanh）}{\kixzsv{6}{392}}
\bkidxline{はかる数}{\kixzsv{2}{19}}
\bkidxline{掃き出し法}{\kixzsv{7}{202}}
\bkidxline{箱ひげ図}{\kixzsv{3}{402}}
\bkidxline{柱の体積}{\kixzsv{4}{169}}
\bkidxline{波数}{\kixzsv{8}{551}}
\bkidxline{パスカルの三角形}{\kixzsv{3}{506}}
\bkidxline{発散}{\kixzsv{8}{350}}
\bkidxline{波動方程式}{\kixzsv{8}{548}}
\bkidxline{ばね定数}{\kixzsv{2}{319}\kixzsvsep \kixzsv{5}{285}}
\bkidxline{速さ}{\kixzsv{2}{65}\kixzsvsep \kixzsv{5}{9}}
\bkidxline{半区間の展開}{\kixzsv{9}{102}}
\bkidxline{半径}{\kixzsv{1}{327}}
\bkidxline{半減期}{\kixzsv{3}{339}\kixzsvsep \kixzsv{5}{353}}
\bkidxline{番地}{\kixzsv{2}{388}}
\bkidxline{反比例}{\kixzsv{2}{321}}
\bkidxline{判別式}{\kixzsv{3}{306}}
\bkidxline{比}{\kixzsv{2}{106}}
\bkidxline{非圧縮}{\kixzsv{8}{369}}
\bkidxline{PID 制御}{\kixzsv{9}{548}}
\bkidxline{ひき算}{\kixzsv{1}{82}}
\bkidxline{－（ひく）}{\kixzsv{1}{82}}
\bkidxline{ヒストグラム}{\kixzsv{3}{396}}
\bkidxline{ひっ算}{\kixzsv{1}{116}}
\bkidxline{ビット}{\kixzsv{2}{368}}
\bkidxline{1つ分}{\kixzsv{1}{204}}
\bkidxline{1目もり}{\kixzsv{1}{101, 290}}
\bkidxline{ひとり分}{\kixzsv{1}{236}}
\bkidxline{比の値}{\kixzsv{2}{106}}
\bkidxline{微分}{\kixzsv{5}{67}}
\bkidxline{微分係数}{\kixzsv{5}{66}}
\bkidxline{微分積分学の基本定理}{\kixzsv{5}{267}}
\bkidxline{微分方程式}{\kixzsv{5}{14}\kixzsvsep \kixzsv{6}{18, 32}}
\bkidxline{百のくらい}{\kixzsv{1}{40}}
\bkidxline{百分率}{\kixzsv{2}{109}}
\bkidxline{表}{\kixzsv{1}{286}}
\bkidxline{秒}{\kixzsv{1}{174}}
\bkidxline{表計算ソフト}{\kixzsv{2}{387}}
\bkidxline{標準偏差}{\kixzsv{3}{421}}
\bkidxline{標本調査}{\kixzsv{3}{512}}
\bkidxline{比例}{\kixzsv{2}{316}}
\bkidxline{フィードバック}{\kixzsv{9}{361}}
\bkidxline{$v$-$t$ グラフ}{\kixzsv{2}{67}}
\bkidxline{フーリエ級数}{\kixzsv{9}{51}}
\bkidxline{フーリエの法則}{\kixzsv{6}{112}}
\bkidxline{フーリエ変換}{\kixzsv{9}{169}}
\bkidxline{フェーザ}{\kixzsv{6}{360}}
\bkidxline{ふえる}{\kixzsv{1}{68}}
\bkidxline{複素数}{\kixzsv{6}{287}}
\bkidxline{複素フーリエ係数}{\kixzsv{9}{131}}
\bkidxline{複素平面}{\kixzsv{6}{315}}
\bkidxline{複利}{\kixzsv{5}{336}}
\bkidxline{符号}{\kixzsv{2}{47}}
\bkidxline{物質微分}{\kixzsv{8}{553}}
\bkidxline{物理シリーズ}{\kixzsv{9}{578}}
\bkidxline{物理定数}{\kixzsv{3}{30}}
\bkidxline{不定積分}{\kixzsv{5}{269}}
\bkidxline{不等号}{\kixzsv{2}{41}}
\bkidxline{負の数}{\kixzsv{2}{40}}
\bkidxline{負の周波数}{\kixzsv{9}{128}}
\bkidxline{部分積分}{\kixzsv{5}{402}}
\bkidxline{部分分数分解}{\kixzsv{9}{322}}
\bkidxline{振り子の周期}{\kixzsv{2}{173}\kixzsvsep \kixzsv{5}{430}}
\bkidxline{プログラム}{\kixzsv{2}{406}}
\bkidxline{ブロック線図}{\kixzsv{3}{101}\kixzsvsep \kixzsv{9}{371}}
\bkidxline{分}{\kixzsv{1}{173}}
\bkidxline{分解能}{\kixzsv{2}{370}}
\bkidxline{分散}{\kixzsv{3}{421}}
\bkidxline{分数}{\kixzsv{1}{270}\kixzsvsep \kixzsv{2}{78}}
\bkidxline{分数の指数}{\kixzsv{3}{342}}
\bkidxline{分度器}{\kixzsv{1}{344}}
\bkidxline{分配法則}{\kixzsv{2}{239}}
\bkidxline{分布定数回路}{\kixzsv{7}{439}}
\bkidxline{分母}{\kixzsv{5}{63}}
\bkidxline{ペア}{\kixzsv{1}{20}}
\bkidxline{平均}{\kixzsv{2}{126}}
\bkidxline{平均の速さ}{\kixzsv{5}{21}}
\bkidxline{平均変化率}{\kixzsv{5}{66}}
\bkidxline{平行移動}{\kixzsv{4}{357}}
\bkidxline{平行四辺形の面積}{\kixzsv{4}{163}}
\bkidxline{平方完成}{\kixzsv{3}{237}}
\bkidxline{平方根}{\kixzsv{2}{172}}
\bkidxline{平方センチメートル（cm\textsuperscript {2}）}{\kixzsv{1}{361}}
\bkidxline{平方メートル（m\textsuperscript {2}）}{\kixzsv{1}{362}}
\bkidxline{ページランク}{\kixzsv{7}{443}}
\bkidxline{べき乗法}{\kixzsv{7}{325}}
\bkidxline{ベクトル}{\kixzsv{4}{344}}
\bkidxline{ベクトルの和}{\kixzsv{4}{347}}
\bkidxline{ベクトル場}{\kixzsv{8}{11, 179}}
\bkidxline{ヘルムホルツの分解}{\kixzsv{8}{582}}
\bkidxline{辺}{\kixzsv{1}{324}}
\bkidxline{変位}{\kixzsv{5}{208}}
\bkidxline{変域}{\kixzsv{3}{31, 89}}
\bkidxline{偏角}{\kixzsv{6}{318}}
\bkidxline{変化の割合}{\kixzsv{3}{136}}
\bkidxline{変化量}{\kixzsv{2}{64}\kixzsvsep \kixzsv{5}{20}}
\bkidxline{変換}{\kixzsv{7}{72}}
\bkidxline{変換表}{\kixzsv{9}{215}}
\bkidxline{偏差}{\kixzsv{3}{418}}
\bkidxline{変数}{\kixzsv{2}{408}\kixzsvsep \kixzsv{3}{27}}
\bkidxline{変数分離}{\kixzsv{6}{140}}
\bkidxline{変数分離法}{\kixzsv{9}{59}}
\bkidxline{偏導関数}{\kixzsv{8}{78}}
\bkidxline{偏微分}{\kixzsv{8}{76}}
\bkidxline{ポアソン分布}{\kixzsv{6}{157}}
\bkidxline{ポアソン方程式}{\kixzsv{8}{506}}
\bkidxline{ボイルの法則}{\kixzsv{3}{202}}
\bkidxline{ぼうグラフ}{\kixzsv{1}{288}}
\bkidxline{方向微分}{\kixzsv{8}{147}}
\bkidxline{方向余弦}{\kixzsv{4}{448}}
\bkidxline{法線ベクトル}{\kixzsv{4}{449}}
\bkidxline{方程式}{\kixzsv{2}{338}}
\bkidxline{放物線}{\kixzsv{3}{261}}
\bkidxline{放物線の軸}{\kixzsv{3}{262}}
\bkidxline{保存力}{\kixzsv{8}{269}}
\bkidxgroup{ま}
\bkidxline{マクスウェル方程式}{\kixzsv{8}{570}}
\bkidxline{マルコフ連鎖}{\kixzsv{7}{393}}
\bkidxline{右上がり}{\kixzsv{1}{309}}
\bkidxline{右手系}{\kixzsv{4}{439}}
\bkidxline{道のり}{\kixzsv{5}{208}}
\bkidxline{ミリメートル（mm）}{\kixzsv{1}{135}}
\bkidxline{ミリリットル（mL）}{\kixzsv{1}{153}}
\bkidxline{無理数}{\kixzsv{2}{188}}
\bkidxline{メートル（m）}{\kixzsv{1}{138}}
\bkidxline{目分量}{\kixzsv{2}{124}}
\bkidxline{目もり}{\kixzsv{1}{53, 96, 186}}
\bkidxline{面積}{\kixzsv{1}{361}}
\bkidxline{面積分}{\kixzsv{8}{303}}
\bkidxline{面密度}{\kixzsv{8}{215}}
\bkidxline{文字式}{\kixzsv{2}{233}}
\bkidxline{ものさし}{\kixzsv{1}{55}}
\bkidxgroup{や}
\bkidxline{約分}{\kixzsv{2}{80}}
\bkidxline{ヤコビアン}{\kixzsv{8}{117}}
\bkidxline{ヤコビ行列}{\kixzsv{7}{84}}
\bkidxline{有効数字}{\kixzsv{2}{129}}
\bkidxline{有理化}{\kixzsv{3}{304}}
\bkidxline{有理数}{\kixzsv{2}{188}}
\bkidxline{余因子展開}{\kixzsv{7}{156}}
\bkidxline{余弦定理}{\kixzsv{4}{241}}
\bkidxline{四端子定数}{\kixzsv{7}{117}}
\bkidxgroup{ら}
\bkidxline{ライプニッツ}{\kixzsv{5}{76}}
\bkidxline{ラジアン}{\kixzsv{4}{64}\kixzsvsep \kixzsv{5}{363}}
\bkidxline{ラプラシアン}{\kixzsv{8}{503}}
\bkidxline{ラプラス変換}{\kixzsv{9}{220}}
\bkidxline{ラプラス方程式}{\kixzsv{8}{505}}
\bkidxline{力線}{\kixzsv{8}{181}}
\bkidxline{力率}{\kixzsv{4}{323}}
\bkidxline{離散コサイン変換}{\kixzsv{9}{542}}
\bkidxline{離散フーリエ変換（DFT）}{\kixzsv{9}{185}}
\bkidxline{立体角}{\kixzsv{8}{320}}
\bkidxline{リットル（L）}{\kixzsv{1}{152}}
\bkidxline{$\lim $（リミット）}{\kixzsv{5}{62}}
\bkidxline{流線}{\kixzsv{8}{181}}
\bkidxline{流束}{\kixzsv{8}{298}}
\bkidxline{流量}{\kixzsv{5}{291}}
\bkidxline{両側スペクトル}{\kixzsv{9}{132}}
\bkidxline{臨界減衰}{\kixzsv{6}{436}}
\bkidxline{累次積分}{\kixzsv{8}{214}}
\bkidxline{累乗}{\kixzsv{2}{152}}
\bkidxline{$\sqrt {\phantom {2}}$（ルート）}{\kixzsv{2}{172}}
\bkidxline{0（れい）}{\kixzsv{1}{54}}
\bkidxline{零点}{\kixzsv{9}{395}}
\bkidxline{列（表の）}{\kixzsv{1}{220}}
\bkidxline{列ベクトル}{\kixzsv{7}{29}}
\bkidxline{連鎖律}{\kixzsv{8}{111}}
\bkidxline{連続の式}{\kixzsv{8}{369}}
\bkidxline{連立微分方程式}{\kixzsv{6}{496}}
\bkidxline{連立方程式}{\kixzsv{3}{165}}
\bkidxline{ロール・ピッチ・ヨー}{\kixzsv{7}{278}}
\bkidxline{ローレンツ力}{\kixzsv{4}{483}}
\bkidxgroup{わ}
\bkidxline{わき出し}{\kixzsv{8}{348}}
\bkidxline{割合}{\kixzsv{2}{108}}
\bkidxline{わり算}{\kixzsv{1}{236}}
\bkidxline{÷（わる）}{\kixzsv{1}{236}}
\bkidxline{わる数}{\kixzsv{1}{253}}
}{（tools/mkseriesidx.py で作る）}
\end{multicols}

%% ---------------- 奥付 ----------------
\bkpadfour
\global\bk@decofalse\bkfolio{false}
\null\vfill
\noindent\begin{minipage}{\linewidth}\small
{\sffamily\bfseries\Large\color{bkDeep}物理を理解するための数学　第9巻}\par\vspace{1mm}
{\sffamily\bfseries\large 波と変化を読みとく}\par\vspace{3mm}
{\color{bkVol}\hrule height 1pt}\vspace{3mm}
\begin{tabular}{@{}l@{\hspace{6mm}}l@{}}
編者 & 「物理を理解するための数学」編集委員会 \\
 & 委員長　小学校長（まとめ役） \\
 & 委員　数学・電気工学・機械工学・教育学・量子力学・建築工学・\\
 & \phantom{委員　}半導体工学の大学教授、幼稚園教諭、塾講師、\\
 & \phantom{委員　}特別支援学校教諭、心理士、科学解説者、小学校教諭 \\
デザイン & ユニバーサルデザイン担当委員 \\
挿絵 & イラスト担当委員 \\
版 & 第1版（第15回編集委員会）　2026年10月 \\
\end{tabular}\par\vspace{3mm}
{\color{bkVol}\hrule height 1pt}\vspace{2mm}
{\footnotesize この本の問題の答えは、巻末の「こたえ」にあります。くわしい解き方は、別冊の問題集と指導書にのせます。}
\end{minipage}
\makeatother
